\documentclass[11pt]{article}
\usepackage[margin=1in]{geometry}
\usepackage[T1]{fontenc}\usepackage[utf8]{inputenc}
\usepackage{lmodern,amsmath,amssymb,amsthm,mathtools,booktabs,microtype}
\usepackage[colorlinks=true,linkcolor=blue,citecolor=blue,urlcolor=blue]{hyperref}
\hypersetup{pdftitle={Bayesian neural networks adapt to approximation error},pdfauthor={ChatGPT-6 Astra; Claude Opus 5.5; Kaarel Hänni; Dmitry Vaintrob}}
\newtheorem{theorem}{Theorem}[section]\newtheorem{lemma}[theorem]{Lemma}
\newtheorem{corollary}[theorem]{Corollary}\newtheorem{proposition}[theorem]{Proposition}
\theoremstyle{definition}\newtheorem{definition}[theorem]{Definition}
\theoremstyle{remark}\newtheorem{remark}[theorem]{Remark}
\newcommand{\E}{\mathbb E}\newcommand{\Prob}{\mathbb P}\newcommand{\R}{\mathbb R}\newcommand{\N}{\mathbb N}
\newcommand{\KL}{\operatorname{KL}}\newcommand{\Rad}{\operatorname{Rad}}\newcommand{\Lip}{\operatorname{Lip}}
\newcommand{\clip}{\operatorname{clip}}\newcommand{\one}{\mathbf1}
\DeclareMathOperator*{\argmin}{arg\,min}
\title{Bayesian neural networks adapt to approximation error}
\author{ChatGPT-6 Astra \and Claude Opus 5.5 \and Kaarel H\"anni \and Dmitry Vaintrob\thanks{Authors are listed in alphabetical order.}}
\date{18 November 2026}
\begin{document}\maketitle
\begin{abstract}
We give finite-posterior approximation guarantees for a dense Bayesian neural network without knowing the width, coefficient bound, approximation error, or noise level of a useful comparison network. The learner has one fixed architecture and independent Gaussian parameter coordinates. Its prediction is the ordinary posterior mean of its raw real output. An empirical free-energy equation selects the Gaussian observation variance. For every prior exponent $\gamma\ge2$, a comparison network of squared risk $R$ and implementation budget $A$ gives excess, up to logarithms, $A R^{1-1/r}n^{-1/r}+A^r/n$. One may take $r$ equal to hidden depth for tanh with at least two hidden layers, and total affine depth for every fixed globally Lipschitz activation. Shallow tanh has the sharper bound $\sqrt{AR/n}+A/n$. A second proof covers every $\gamma\ge1$ and every fixed globally Lipschitz activation, with a dimensional cost and a noiseless rate $n^{-2/(d+2)}$. All guarantees are independent of learner width and simultaneous over eligible comparison caps. When two calibrations apply, the learner can select the posterior with the smaller observable risk certificate. No architecture prior or mixture of predictors is introduced. A shallow ReLU specialization has sample cost quadratic in path variation, and a normalization argument extends the raw-output bounds to sub-Gaussian labels with only logarithmic changes. We also prove that fixed physical observation variance, and several ordinary variance-hyperprior alternatives, cannot provide this all-width behavior even for a constant target.
\end{abstract}

\section{Learning at an unknown approximation scale}
A small neural network may approximate a regression function well even when no small network explains the labels exactly. A useful Bayesian guarantee should compare test loss with that approximation, remain valid when the learner is much wider, and improve when the comparison error is small. It should also select its likelihood without being told which width or coefficient cap gives a useful approximation.

This paper obtains those properties by retaining empirical free energy in the posterior-tail calculation. The free energy measures the compromise between fit and prior cost that the posterior actually makes on the observed sample. A scalar calibration adapts the likelihood to this compromise. Prediction uses one posterior at one selected variance; neither network size nor temperature is averaged over a prior.

Let $R(g)=\E(g(X)-Y)^2$, let $H$ denote hidden depth, and put $L=H+1$ for depth including the linear readout. For a width-$m$, coefficient-$B$ teacher, useful implementation budgets are
\begin{align}
 E_{\tanh}&=d(d+1)B^2+2Hm^2B^2,\label{eq:intro-tanh-energy}\\
 E_{\mathrm{ReLU}}&=L[d(d+1)]^{1/L}m^{2H/L}B^2.\label{eq:intro-relu-energy}
\end{align}
The first budget also applies to a general Lipschitz activation. Precise costs below add explicit logarithmic neighborhood and confidence terms. At fixed dimension, depth, and activation, write $A_g=\widetilde O(E_g)$ for those costs.

For prior exponents $\gamma\ge2$, the main approximation bound is
\begin{equation}
 R(\widehat f_n)\le\inf_g\left\{R(g)+\widetilde O\left(
 A_gR(g)^{1-1/r}n^{-1/r}+A_g^r/n\right)\right\}.
 \label{eq:intro-oracle}
\end{equation}
Here $r=H$ for tanh with $H\ge2$, and $r=L$ for ReLU. The latter exponent applies to every fixed globally Lipschitz activation, including tanh, GELU, sigmoid and softplus, and allows nonzero activation values at the origin. A scalar-link entropy argument adds only logarithms in sample size; a simpler signed-hinge decomposition avoids those extra logarithms for activations with bounded derivative variation. With one hidden layer, the Gaussian moment hypothesis holds for all $\gamma\ge1$. Shallow tanh has the stronger remainder
\begin{equation}
 \widetilde O\left(\sqrt{A_gR(g)/n}+A_g/n\right).
 \label{eq:intro-shallow}
\end{equation}

For every $\gamma\ge1$, including deep mean-field networks, an input-domain covering argument gives
\begin{equation}
 R(\widehat f_n)\le\inf_g\left\{R(g)+\widetilde O\left(
 A_gR(g)^{1-1/r}n^{-2/[r(d+2)]}
       +A_g^r n^{-2/(d+2)}\right)\right\}.
 \label{eq:intro-all-gamma}
\end{equation}
Every fixed globally Lipschitz activation permits $r=L$. For tanh this improves to
$r=\max\{2,(2+dL)/(d+2)\}$. No bounded-variation assumption on the derivative is needed in this second argument. The dimensional exponent is its cost for covering all prior scales.

The infima have a precise simultaneous interpretation: one sample event compares the selected learner with the population minimizer of every integer width and coefficient cap fitting inside its architecture. Neither those caps nor their risks enter the calibration. The data law may be discrete or continuous and need not follow a Gaussian noise model. For a conditional mean $f_0(x)=\E[Y\mid X=x]$ and noise variance $\nu=\E\operatorname{Var}(Y\mid X)$, the identity
$R(g)=\nu+\|g-f_0\|_{L^2(P_X)}^2$ turns each statement into an approximation guarantee for $f_0$.

The main restriction on comparison functions is a coefficient bound, in addition to width and depth. These raw-output theorems allow arbitrary input laws. The sample-polynomial degree still worsens with depth; this paper does not assert a polynomial bound jointly in width and unrestricted depth. It also does not give an efficient method for evaluating the evidence or posterior integral.

\section{One Gaussian architecture and its likelihood calibration}
Inputs lie in $[-1,1]^d$, labels satisfy $|Y|\le1$, and observations are IID. Appendix~\ref{sec:subgaussian-adaptive-transfer} extends the raw-output guarantees to sub-Gaussian labels; Appendix~\ref{sec:shallow-ReLU-path} specializes the shallow ReLU case to Euclidean-ball inputs. At learner width $N$, put $\bar x=(1,x)$ and
\[
 h^1=\phi(W^1\bar x),\qquad h^\ell=\phi(W^\ell h^{\ell-1})\ (2\le\ell\le H),
 \qquad z_w=v^{\mathsf T}h^H.
\]
There are no internal biases. Ordinary biased teachers can be compiled by a constant neuron, as explained below. The activation is fixed and known. Define
\begin{equation}
 Q_N(w)=\frac dN\|W^1\|_F^2+\sum_{\ell=2}^H\|W^\ell\|_F^2+N\|v\|_2^2.
 \label{eq:energy}
\end{equation}
For $s=N^{-\gamma}$, $\gamma\ge1$, the prior $\mu_{N,s}$ has independent centered Gaussian coordinates with variances $Ns/d$, $s$, and $s/N$ in the first, internal, and readout layers. Its density is proportional to $e^{-Q_N/(2s)}$. These are actual finite-width priors. The full-energy argument requires $\gamma\ge2$ at more than one hidden layer; the operator-norm argument covers the entire specified range.

For a sample $S=((X_i,Y_i))_{i=1}^n$, write
$\widehat R_S(f)=n^{-1}\sum_i(f(X_i)-Y_i)^2$.
At a likelihood parameter $\lambda>0$, choose Gaussian observation variance $n\lambda s$. The posterior and empirical free energy are
\begin{align}
 \Pi_{S,\lambda}(dw)&=Z_S(\lambda)^{-1}
 e^{-\widehat R_S(z_w)/(2\lambda s)}\mu_{N,s}(dw),\label{eq:posterior}\\
 F_S(\lambda)&=-2\lambda s\log Z_S(\lambda)
 =\inf_\rho\left\{\int\widehat R_S(z_w)d\rho
                  +2\lambda s\KL(\rho\Vert\mu_{N,s})\right\}.
 \label{eq:freeenergy}
\end{align}
The Gaussian likelihood's weight-independent normalizer is omitted from $Z_S$. This does not change the posterior at a fixed $\lambda$, but it is part of the definition of the calibration. The procedure is not maximum marginal-likelihood estimation of a physical noise variance.

For the full-energy proof with $r\ge2$, choose the unique positive root
\begin{equation}
 \lambda^r=256a_n^2[2F_S(\lambda)+\lambda(K_n+1)]^{r-1}.
 \label{eq:root-calibration}
\end{equation}
The deterministic $a_n=\widetilde O(n^{-1/2})$ and logarithmic $K_n$ are specified in Section~\ref{sec:abstract}. For shallow tanh, choose the smallest minimizer of
\begin{equation}
 F_S(\lambda)(1+4a_n^2/\lambda)+c_0\lambda
 \quad\text{over }\lambda\ge16a_n^2.
 \label{eq:shallow-calibration}
\end{equation}
Both rules are well defined at every $n\ge2$.

The all-scale proof uses the equivalent notation $\omega=1/s$ and $\tau^2=Ns$. Its root is
\[
 \lambda^p=\theta_n[2F_S(\lambda)+\lambda b_n]^q,
 \qquad p=r(d+2)/2,\quad q=p-(d+2)/2,
\]
with explicit constants and a deterministic small-sample fallback in Section~\ref{ag-sec:adaptive-core}. Each calibration has an observable upper certificate for posterior-average test risk. Section~\ref{sec:certificate-selection} explains how to select the smaller certificate when more than one is available. In every case the final predictor is one ordinary raw posterior mean,
\[
 \widehat f_n(x)=\int z_w(x)d\Pi_{S,\lambda_S}(w).
\]

\section{A Gaussian free-energy oracle theorem}
\label{sec:abstract}
We first isolate the statistical argument. It applies to a family of learners, indexed by $i$, whose outputs are $z_{i,w}$ and whose nonnegative energy functions are $Q_i$. Under the prior $\mu_i$, assume $Q_i/s_i$ has a chi-square distribution with $D_i$ degrees of freedom, and
\begin{equation}
 0<s_i\le1,\qquad s_iD_i\le D_0,
 \qquad \E_{\mu_i}\widehat R_S(z_{i,w})\le M.
 \label{eq:abstract-prior}
\end{equation}
The last bound is required for every allowed sample, with the same finite $M$. We will use $M=2$ for both neural-network activations. A full Gaussian quadratic energy satisfies this assumption, as does the energy of a Gaussian coordinate block when other parameter coordinates have an independent, possibly restricted prior. Only the stated energy law, amplitude bound, and capacity bound are used below.

Fix an integer $r\ge1$. Assume the energy-ball union
$\mathcal Z_q=\bigcup_i\{z_{i,w}:Q_i(w)\le q\}$ has a countable determining subclass and, on every list of $n$ inputs,
\begin{equation}
 |z_{i,w}(x)|\le V(1+Q_i(w))^{r/2},\qquad
 \Rad_n(\mathcal Z_q)\le\frac{\kappa_n(1+q)^{r/2}}{\sqrt n}.
 \label{eq:abstract-capacity}
\end{equation}
Here $V\ge1$ and $\kappa_n$ are independent of the learner index. Width unions are therefore allowed.

We use the following smooth-loss inequality~\cite[Theorem~1]{sst}. For a class $\mathcal F$ on which a nonnegative $s_0$-smooth loss is bounded by $b$, with probability at least $1-\eta$, simultaneously for $f\in\mathcal F$,
\begin{align}
 R(f)-\widehat R_S(f)\le C_{\rm sm}\bigg[&\sqrt{\widehat R_S(f)}
 \left(\sqrt{s_0}\log^{3/2}(2n)\Rad_n(\mathcal F)
       +\sqrt{b\log(1/\eta)/n}\right)\notag\\
 &+s_0\log^3(2n)\Rad_n(\mathcal F)^2+b\log(1/\eta)/n\bigg],
 \label{eq:sst}
\end{align}
where $C_{\rm sm}\ge1$ is universal and $\Rad_n$ denotes worst-case empirical Rademacher complexity. For raw square loss $s_0=2$.

Define the deterministic constants
\begin{align}
 B_r&=2^rV^2+2,&M_r&=1+(2D_0+4r)^r,\notag\\
 K_n&=2D_0\log2+4\log(B_rM_rn^4)+1,&
 c_0&=\frac18+\frac{D_0\log2}{4},\label{eq:tail-constants}\\
 J_n&=\left\lceil\log_2(2Mn+K_n+2)\right\rceil,&
 \ell_n&=\max\left\{1,\log\frac{2(J_n+1)}\delta\right\},\notag\\
 a_n&=\max\left\{\frac1{\sqrt n},\,
 2^{r/2}C_{\rm sm}\left[
 \frac{\sqrt2\kappa_n\log^{3/2}(2n)}{\sqrt n}
 +(V+1)\sqrt{\frac{\ell_n}n}\right]\right\}.\label{eq:an}
\end{align}
The confidence parameter here is the failure probability allocated to the common capacity event. Later, half the desired failure probability is allocated to that event.

\begin{lemma}[One statistical event]
\label{lem:stat-event}
With probability at least $1-\delta$, simultaneously for all learner indices and all parameters with $Q_i(w)\le2Mn+K_n+1$,
\begin{equation}
 R(z_{i,w})-\widehat R_S(z_{i,w})
 \le a_n\sqrt{\widehat R_S(z_{i,w})}(1+Q_i(w))^{r/2}
       +a_n^2(1+Q_i(w))^r.
 \label{eq:stat-event}
\end{equation}
\end{lemma}
\begin{proof}
Apply \eqref{eq:sst} to $\mathcal Z_{2^j}$, $0\le j\le J_n$, with failure probability $\delta/(J_n+1)$. The square-loss bound on that class is $(V+1)^2(1+2^j)^r$. For $Q\le1$ use $j=0$; otherwise use the least $j$ with $Q\le2^j$. In either case $1+2^j\le2(1+Q)$. Substitution of \eqref{eq:abstract-capacity} into \eqref{eq:sst} and the definition of $a_n$ give \eqref{eq:stat-event}. The cross term in the square defining $a_n^2$ is nonnegative, so that square dominates the sum of the two quadratic terms. A finite union bound proves the claim.
\end{proof}

For the rest of this section, suppress the learner index and write $s,Q,D,\mu,F,\Pi_\lambda$ for one member of the family. All deterministic estimates remain uniform over its members.

\begin{lemma}[Weighted objective tails]
\label{lem:joint-tail}
For every sample, every $\lambda>0$, and every $u>0$,
\begin{equation}
 \int(1+Q^r)\one_{\{\widehat R_S+\lambda Q>u\}}d\Pi_\lambda
 \le M_r\exp\left[\frac{F(\lambda)}{2\lambda s}
                   -\frac{u}{4\lambda s}+\frac{D\log2}{2}\right].
 \label{eq:joint-tail}
\end{equation}
For $u=2F(\lambda)+\lambda K_n$, the raw-risk integral on this tail is at most $n^{-4}$.
\end{lemma}
\begin{proof}
Put $J=\widehat R_S+\lambda Q$. On $J>u$,
\[
 e^{-\widehat R_S/(2\lambda s)}
 =e^{-J/(2\lambda s)+Q/(2s)}
 \le e^{-u/(4\lambda s)}e^{Q/(4s)},
\]
after dropping the remaining nonnegative empirical loss. The normalizer is $e^{-F/(2\lambda s)}$. Exponentially tilting the chi-square law by $e^{Q/(4s)}$ doubles its scale and contributes the factor $2^{D/2}$. Under that tilted law,
\[
 \E Q^r=(2s)^r\prod_{j=0}^{r-1}(D+2j)
 \le(2D_0+4r)^r.
\]
This proves \eqref{eq:joint-tail}, even when $Q$ describes only a Gaussian block of the full prior. Square loss is at most
$2V^2(1+Q)^r+2\le B_r(1+Q^r)$.
At the specified $u$, the exponent in \eqref{eq:joint-tail} is at most
$-[K_n-2D_0\log2]/(4s)\le-\log(B_rM_rn^4)$.
Multiply by $B_r$ to obtain the raw-risk assertion.
\end{proof}

\begin{lemma}[Relative entropy pays quadratic energy]
\label{lem:entropy}
For every law $\rho$ relative to $\mu$,
\begin{equation}
 \int Qd\rho\le4s\KL(\rho\Vert\mu)+2D_0\log2.
 \label{eq:entropy}
\end{equation}
Also $\int\widehat R_Sd\Pi_\lambda\le F(\lambda)\le M$.
\end{lemma}
\begin{proof}
The exact Gaussian moment is $\E_\mu e^{Q/(4s)}=2^{D/2}$. The entropy inequality gives
$\KL(\rho\Vert\mu)\ge\int Qd\rho/(4s)-(D/2)\log2$.
The remaining statements follow from the variational identity for $F$, nonnegativity of relative entropy, and the choice $\rho=\mu$.
\end{proof}

\begin{lemma}[The root is unique and measurable]
\label{lem:root}
For $r\ge2$, equation \eqref{eq:root-calibration} has one positive solution. It obeys
\begin{equation}
 \lambda_S\ge256a_n^2(K_n+1)^{r-1}\ge1/n.
 \label{eq:root-floor}
\end{equation}
The solution is a measurable function of the sample; for continuous network outputs it is continuous.
\end{lemma}
\begin{proof}
The variational formula makes $F$ finite, increasing, and concave on $(0,\infty)$, hence continuous. Directly from its Laplace representation, $F(\lambda)/\lambda$ is decreasing. Therefore
\[
 \Psi(\lambda)=\frac{\lambda^r}{[2F(\lambda)+\lambda(K_n+1)]^{r-1}}
 =\frac{\lambda}{[2F(\lambda)/\lambda+K_n+1]^{r-1}}
\]
is continuous and strictly increasing. Its limit at zero is zero, since it is at most $\lambda/(K_n+1)^{r-1}$. Its limit at infinity is infinity, since $F\le M$. Thus it crosses $256a_n^2$ exactly once. Its upper bound at small $\lambda$ gives \eqref{eq:root-floor}.
For each fixed $t>0$, $\{\lambda_S<t\}=\{\Psi_S(t)>256a_n^2\}$ is measurable. If empirical loss is continuous in the sample, bounded convergence makes the Laplace integral jointly continuous in sample and positive $\lambda$. Continuity of the strictly monotone root follows from its crossing property.
\end{proof}

\begin{theorem}[Deep free-energy calibration]
\label{thm:deep-abstract}
Let $r\ge2$. On the event of Lemma~\ref{lem:stat-event}, simultaneously for every learner in the family, the posterior with calibration \eqref{eq:root-calibration} satisfies
\begin{equation}
 \int R(z_w)d\Pi_{\lambda_S}\le F(\lambda_S)+c_0\lambda_S+n^{-4}.
 \label{eq:freeenergy-certificate}
\end{equation}
The ordinary posterior mean has the same risk bound.
Suppose, on this sample, a comparison gives
$F(\lambda_S)\le u+\lambda_S A$, with $u,A\ge0$, and put
$\mathcal B=2A+K_n+1+c_0$. Then for a constant $C_r$ depending only on $r$,
\begin{equation}
 \int R(z_w)d\Pi_{\lambda_S}
 \le u+C_r\left[\mathcal B u^{1-1/r}a_n^{2/r}
                         +\mathcal B^r a_n^2\right]+n^{-4}.
 \label{eq:deep-empirical-oracle}
\end{equation}
All valid comparisons may be made on this same event.
\end{theorem}
\begin{proof}
Write $\lambda=\lambda_S$ and $G=2F(\lambda)+\lambda(K_n+1)$. On the complement of the tail in Lemma~\ref{lem:joint-tail},
\[
 \widehat R_S\le G,\qquad 1+Q\le G/\lambda,
 \qquad Q\le2M/\lambda+K_n\le2Mn+K_n.
\]
The last inequality uses \eqref{eq:root-floor}; thus the fixed statistical event applies. Since $r\ge2$, its two terms obey
\begin{align*}
 a_n\sqrt{\widehat R_S}(1+Q)^{r/2}
 &\le\lambda(1+Q)\frac{a_nG^{(r-1)/2}}{\lambda^{r/2}}
 =\frac{\lambda(1+Q)}{16},\\
 a_n^2(1+Q)^r
 &\le\lambda(1+Q)\frac{a_n^2G^{r-1}}{\lambda^r}
 =\frac{\lambda(1+Q)}{256}.
\end{align*}
Their sum is at most $\lambda(1+Q)/8$. Integrate, extend the nonnegative empirical-loss and energy terms to the whole parameter space, and add the weighted tail. Lemma~\ref{lem:entropy} gives
\begin{align*}
 \int Rd\Pi_\lambda
 &\le\int\widehat R_Sd\Pi_\lambda+\frac\lambda8\int Qd\Pi_\lambda
                  +\frac\lambda8+n^{-4}\\
 &\le F(\lambda)+\lambda\left(\frac18+\frac{D_0\log2}{4}\right)+n^{-4}.
\end{align*}
The posterior has finite raw second moments, so Jensen gives its mean prediction.

For the comparison, set $p=(r-1)/r$ and $v=256a_n^2$. The root and the assumed free-energy bound give
\[
 \lambda\le v^{1/r}[2u+(2A+K_n+1)\lambda]^p
 \le v^{1/r}(2u)^p+v^{1/r}(2A+K_n+1)^p\lambda^p.
\]
If the first term is at least $\lambda/2$, it bounds $\lambda$ directly. Otherwise the second is at least $\lambda/2$, yielding
$\lambda\le2^rv(2A+K_n+1)^{r-1}$. Combining the cases gives
\[
 \lambda\le C_r\left[a_n^{2/r}u^{1-1/r}
                  +a_n^2(2A+K_n+1)^{r-1}\right].
\]
Substitute this in \eqref{eq:freeenergy-certificate} and use
$A+c_0\le\mathcal B$ to obtain \eqref{eq:deep-empirical-oracle}.
\end{proof}

\begin{theorem}[Calibration for square-root norm capacity]
\label{thm:shallow-abstract}
Let $r=1$. On the event of Lemma~\ref{lem:stat-event}, simultaneously for every learner and every $\lambda\ge16a_n^2$,
\begin{equation}
 \int R(z_w)d\Pi_\lambda
 \le F(\lambda)(1+4a_n^2/\lambda)+c_0\lambda+n^{-4}.
 \label{eq:shallow-certificate}
\end{equation}
Choose the smallest minimizer of the right side without $n^{-4}$, as in \eqref{eq:shallow-calibration}. This selection is measurable. If a comparison satisfies
$F(\lambda)\le u+\lambda A$ throughout the interval, where $u,A\ge0$, and $\mathcal B=A+c_0$, then the selected posterior obeys
\begin{equation}
 \int R(z_w)d\Pi_{\lambda_S}
 \le u+4a_n\sqrt{\mathcal B u}+20a_n^2\mathcal B+n^{-4}.
 \label{eq:shallow-empirical}
\end{equation}
Again the same bound holds for its ordinary raw posterior mean, and all comparisons can be made on the same event.
\end{theorem}
\begin{proof}
The interval has $\lambda\ge16/n\ge1/n$, so the good-objective region is covered by the same statistical event. Write $U=\int\widehat R_Sd\Pi_\lambda$ and $T=1+\int Qd\Pi_\lambda$. Integrating \eqref{eq:stat-event} on that region and using Cauchy--Schwarz bounds its two contributions by
\[
 a_n\sqrt{UT}+a_n^2T
 \le\frac\lambda8T+\frac{4a_n^2}{\lambda}U.
\]
Here Young's inequality gives the first $\lambda T/16$, and $\lambda\ge16a_n^2$ gives the second. Apply the entropy inequality to $\lambda\int Q/8$, use $U\le F$, and add the raw-risk tail to obtain \eqref{eq:shallow-certificate}.

The certificate is continuous and at least $c_0\lambda$, while at $\lambda_0=16a_n^2$ it is at most $5M/4+c_0\lambda_0$. Every minimizer therefore lies in the deterministic compact interval $[\lambda_0,\lambda_0+5M/(4c_0)]$. It attains a minimum there. For every fixed $t$ in this interval, the event that its smallest minimizer is at most $t$ is the equality of the minimum over the interval ending at $t$ and the minimum over the full interval. These minima are measurable, since they equal infima over countable dense sets. This proves measurability.

For any allowed $\lambda$, the teacher comparison bounds the certificate by
\[
 u+\mathcal B\lambda+4a_n^2u/\lambda+4a_n^2A.
\]
Choose $\lambda=\max\{16a_n^2,2a_n\sqrt{u/\mathcal B}\}$, interpreting the square-root term as zero when $u=0$. The bound
$\lambda\le16a_n^2+2a_n\sqrt{u/\mathcal B}$ controls the linear term, while
$4a_n^2u/\lambda\le2a_n\sqrt{\mathcal B u}$. This proves \eqref{eq:shallow-empirical}. Jensen gives the mean-prediction statement.
\end{proof}

\section{The neural-network ingredients}
\label{sec:network-ingredients}
We verify the hypotheses of the abstract theorem, including the comparison cost, without referring to a separate manuscript.

\subsection{Gaussian moments and comparison neighborhoods}
The parameter dimension in \eqref{eq:energy} is
\[
 D_N=N(d+2)+(H-1)N^2.
\]
When $s=N^{-\gamma}$ and $\gamma\ge2$, or when $H=1$ and $\gamma\ge1$, one has
\begin{equation}
 Ns\le1,\qquad D_Ns\le D_0:=d+H+1.
 \label{eq:dimension-speed}
\end{equation}
Under the prior, $Q_N=s\chi^2_{D_N}$. The independent centered readout makes the prior prediction mean zero. For tanh, its conditional variance is $(s/N)\|h^H\|^2\le s$, hence its prior expected empirical loss is at most two. For ReLU, conditional Gaussian symmetry gives
\[
 \E\frac{\|h^1\|^2}N\le Ns,\qquad
 \E\frac{\|h^\ell\|^2}N=\frac{Ns}{2}\E\frac{\|h^{\ell-1}\|^2}N.
\]
The readout multiplies the last normalized second moment by $s$. Thus its prior output second moment is at most $s$, and again $M=2$ is valid.

\begin{lemma}[Exact implementations at every larger width]
\label{lem:teacher-implementation}
A width-$m$ teacher of depth $L=H+1$ and coefficient cap $B\ge1$ has, at every learner width $N\ge m$, an exact implementation of energy at most $E_\phi$ in \eqref{eq:intro-tanh-energy} or \eqref{eq:intro-relu-energy}, respectively. Its output is bounded by
\begin{equation}
 U_{\tanh}=mB,\qquad U_{\mathrm{ReLU}}=(d+1)m^HB^L.
 \label{eq:teacher-amplitudes}
\end{equation}
\end{lemma}
\begin{proof}
Partition each learner hidden layer into $m$ groups of sizes
$k_j\in\{\lfloor N/m\rfloor,\lceil N/m\rceil\}$. Copy each teacher first-layer row into its group. Replace a teacher edge leaving group $j$ by the original weight divided by $k_j$, on every corresponding edge. Divide each readout weight similarly. Every clone has its teacher activation, so the output is unchanged. The first weighted layer energy is at most $d(d+1)B^2$. Each internal layer and the readout have energy at most $2m^2B^2$, since group-size ratios are at most two and $k_j\ge N/(2m)$. This proves the tanh budget.

For ReLU, first zero-pad the teacher into the learner architecture, multiply its first matrix by $\sqrt N$, and divide its readout by $\sqrt N$. Homogeneity preserves the exact function. The weighted layer energies are now bounded by
\[
 d m(d+1)B^2,\quad m^2B^2\text{ in each internal layer},\quad mB^2.
\]
Their product is at most $d(d+1)m^{2H}B^{2L}$. A zero layer gives the zero function. Otherwise multiply each layer by a positive factor, with product one, to make all weighted layer energies equal to their geometric mean. This again preserves the function and gives \eqref{eq:intro-relu-energy}.
The amplitude estimates follow by direct propagation through the original teacher: tanh features are bounded by one, while ReLU contributes at most a factor $mB$ at each layer after the first, whose amplitude is at most $(d+1)B$.
\end{proof}

\begin{lemma}[A Gaussian teacher neighborhood]
\label{lem:teacher-neighborhood}
Let $g$ have an exact implementation of energy at most $E$ and amplitude at most $U$. Set
\begin{equation}
 K_g=\sqrt2L(\sqrt E+1)^H,
 \qquad A_{g,n}=E+1+2D_0\log[4K_g(U+1)n].
 \label{eq:teacher-cost}
\end{equation}
For every sample, every eligible width and scale in \eqref{eq:dimension-speed}, and every $\lambda\ge1/n$,
\begin{equation}
 F_S(\lambda)\le\widehat R_S(g)+\lambda A_{g,n}.
 \label{eq:teacher-freeenergy}
\end{equation}
\end{lemma}
\begin{proof}
Let $w_*$ denote the exact implementation and perturb it by an independent Gaussian with covariance $\varepsilon^2$ times the prior covariance, $0<\varepsilon\le1$. The activations tanh and ReLU are both one-Lipschitz and vanish at zero. Normalize hidden Euclidean norms by $\sqrt N$. Every deterministic internal operator norm, and the normalized readout norm, is at most $\sqrt E$. If $b_\ell$ is the $L^2$ normalized size of the perturbed hidden layer, independence of the current Gaussian perturbation gives
\[
 b_1\le\sqrt2(\sqrt E+\varepsilon),\qquad
 b_\ell\le(\sqrt E+\varepsilon)b_{\ell-1},
 \qquad b_\ell\le\sqrt2(\sqrt E+\varepsilon)^\ell.
\]
For the normalized $L^2$ difference $\Delta_\ell$ from the deterministic state,
\[
 \Delta_1\le\sqrt2\varepsilon,\qquad
 \Delta_\ell\le\sqrt E\Delta_{\ell-1}+\varepsilon b_{\ell-1}
 \le\sqrt2\ell\varepsilon(\sqrt E+\varepsilon)^{\ell-1}.
\]
The readout therefore gives pointwise output root mean square error at most
$\sqrt E\Delta_H+\varepsilon b_H\le\varepsilon K_g$.
Only $Ns\le1$ and $s\le1$ were used in the noise bounds.

Choose the fixed perturbation size $\varepsilon=1/[4K_g(U+1)n]$. Minkowski, or direct expansion of square loss, gives expected empirical loss at most $\widehat R_S(g)+1/n\le\widehat R_S(g)+\lambda$. For the comparison Gaussian $\rho$, the exact Gaussian relative-entropy formula gives
\[
 2\lambda s\KL(\rho\Vert\mu)
 \le\lambda E+2\lambda D_0\log(1/\varepsilon).
\]
Apply the variational formula \eqref{eq:freeenergy} with this fixed perturbation. This yields \eqref{eq:teacher-freeenergy}.
\end{proof}

\subsection{Tanh capacity}
For this subsection only, omit the first-layer energy and put
\[
 Q_N^+=\sum_{\ell=2}^H\|W^\ell\|_F^2+N\|v\|_2^2,
 \qquad D_n^{\tanh}=\sqrt{2(d+1)\log(n+1)}+\frac{H-1}{2}.
\]

\begin{lemma}[Tanh norm capacity]
\label{lem:tanh-capacity}
On every list of $n$ inputs,
\begin{equation}
 \Rad_n\left(\bigcup_N\{z_w:Q_N^+\le q\}\right)
 \le (q/H)^{H/2}\frac{D_n^{\tanh}}{\sqrt n}.
 \label{eq:tanh-capacity}
\end{equation}
In particular, \eqref{eq:abstract-capacity} holds for the full energy with $r=H$, $V=1$, and $\kappa_n=D_n^{\tanh}/H^{H/2}$, including $r=1$ when $H=1$.
\end{lemma}
\begin{proof}
Fix the input list. Rademacher complexity is unchanged by taking closed convex hulls. For $t\in[0,1]$, write $\clip_t(u)=\max(-t,\min(t,u))$. If $\operatorname{Clip}(A)$ applies a common such threshold to vectors of a bounded set $A\subset\R^n$, the identity
$\clip_t(u)=(|u+t|-|u-t|)/2$, contraction, and symmetry of the Rademacher signs give
\begin{equation}
 \Rad(\operatorname{Clip}(A))\le\Rad(A)+\frac1{2\sqrt n}.
 \label{eq:variable-clip}
\end{equation}
Indeed, split the supremum between the two absolute-value expressions and contract each. Their added terms sum to
$\E|\sum_i\epsilon_i|/(2n)\le1/(2\sqrt n)$.

For $a>0$ and $|u|\le1$ there is the positive mixture identity
\begin{equation}
 \frac{\tanh(au)}a=
 \int_0^1\clip_t(u)\,2a\tanh(at)\operatorname{sech}^2(at)dt
 +\operatorname{sech}^2(a)\clip_1(u).
 \label{eq:tanh-mixture}
\end{equation}
Its weights are nonnegative and sum to one. Integration by parts with $\tanh^2(at)$ verifies the formula for $u\ge0$; oddness gives the rest.

Let $\mathcal F_1$ be the closed convex hull of all first-layer tanh feature vectors, with unrestricted first-layer rows. Recursively take
$\mathcal F_\ell=\overline{\operatorname{conv}}\operatorname{Clip}(\mathcal F_{\ell-1})$. These sets are symmetric, convex, contain zero, and lie in $[-1,1]^n$. With $\rho^1=\boldsymbol1_N$ and
$\rho^\ell=|W^\ell|\rho^{\ell-1}$, induction using \eqref{eq:tanh-mixture} puts neuron $j$'s sample vector in $\rho^\ell_j\mathcal F_\ell$. For a positive scaling, divide its preactivation by that scaling to obtain a convex signed combination of earlier normalized features; a zero scaling forces a zero neuron.
The output sample vector consequently lies in
$P(w)\mathcal F_H$, where
\[
 P(w)=|v|^{\mathsf T}|W^H|\cdots|W^2|\boldsymbol1_N
 \le\sqrt N\|v\|_2\prod_{\ell=2}^H\|W^\ell\|_F
 \le(Q_N^+/H)^{H/2}.
\]
The last inequality is arithmetic--geometric mean on the $H$ energy terms.

A real random variable with distribution function $(1+\tanh t)/2$ gives
$\tanh u=\E\operatorname{sign}(u-T)$. Thus $\mathcal F_1$ lies in the convex hull of affine-halfspace sign vectors on the inputs. There are at most $(n+1)^{d+1}$ such vectors by the halfspace growth bound. The finite-class Rademacher bound gives complexity at most
$\sqrt{2(d+1)\log(n+1)/n}$. Apply \eqref{eq:variable-clip} $H-1$ times, and then multiply by the preceding bound on $P(w)$.
Finally $|z_w|\le\sqrt{Q_N^+}\le(1+Q_N)^{H/2}$. This verifies the stated abstract parameters. All suprema have countable determining subclasses: rational parameter approximations with strict energy inequalities are dense, and the width union is countable.
\end{proof}

\subsection{ReLU capacity}
\begin{lemma}[ReLU norm capacity]
\label{lem:relu-capacity}
For $r=L=H+1$, \eqref{eq:abstract-capacity} holds with
\begin{equation}
 V=1,\qquad \kappa_n=2^H\sqrt{\frac{d+1}{d}}L^{-L/2}.
 \label{eq:relu-capacity-constant}
\end{equation}
\end{lemma}
\begin{proof}
Let $P=\|v\|_2\prod_{\ell=1}^H\|W^\ell\|_F$. The weighted first-layer and readout factors cancel $N$, and arithmetic--geometric mean gives
$P\le d^{-1/2}(Q_N/L)^{L/2}$. Lipschitzness and positive homogeneity imply
$|z_w(x)|\le\sqrt{d+1}P\le Q_N^{L/2}$, since
$\sqrt{(d+1)/d}L^{-L/2}\le1$.

Positive homogeneity normalizes all hidden Frobenius norms to one and leaves the readout norm at most $P$. A zero layer gives the zero function. For unit $P$, optimize the readout to replace the signed sum on a fixed sample by the Euclidean norm of the signed hidden-feature sums. Write each final hidden row as its norm times a unit row. Homogeneity and the sum-of-squares Frobenius constraint bound that Euclidean norm by the supremum of a single unit-row signed sum. Absolute Rademacher contraction for ReLU bounds its expectation by twice that of the network with a linear readout and one fewer hidden layer. Repeat this $H$ times. The remaining unit-norm linear class on $\bar x$ has complexity at most $\sqrt{(d+1)/n}$. Multiplication by $P$ yields the stated bound. These optimizations allow the supremum over all widths at each step. Readout sign reversal makes the absolute and ordinary conventions agree for the output class. Continuity and rational parameter approximations provide a countable determining subclass.
\end{proof}

\section{Simultaneous neural-network oracle inequalities}
\label{sec:network-corollaries}
Fix $d,H$, and a desired failure probability $0<\delta<1$. In the constants of Section~\ref{sec:abstract}, use confidence $\delta/2$, $D_0=d+H+1$, $M=2$, and the capacity parameters of Lemma~\ref{lem:tanh-capacity} or Lemma~\ref{lem:relu-capacity}. The calibration is completely specified by these choices and the learner's empirical free energy.

For positive integers $m,b$, let $\mathcal G_{m,b}$ be the comparison networks of hidden width at most $m$, depth $L$, and coefficients bounded in absolute value by $b$. Zero padding represents this as a compact finite-dimensional parameter set. Population square loss is continuous on it, because inputs are bounded and network outputs have a common finite bound. Thus a deterministic minimizer $g_{m,b}$ exists. Write
\[
 R_{m,b}=\min_{g\in\mathcal G_{m,b}}R(g)\le1.
\]
The last inequality holds because the zero network belongs to the class.

Use $E_{m,b},U_{m,b}$ from Lemma~\ref{lem:teacher-implementation}, and define the completely explicit comparison costs
\begin{align}
 K_{m,b}&=\sqrt2L(\sqrt{E_{m,b}}+1)^H,\notag\\
 A_{m,b,n}&=E_{m,b}+1+2D_0\log[4K_{m,b}(U_{m,b}+1)n],\notag\\
 T_{m,b,n}&=2A_{m,b,n}+K_n+1+c_0,\label{eq:simultaneous-cost}\\
 \ell_{m,b}&=\log\frac{32(m+1)^2(b+1)^2}\delta,\qquad
 \mathcal A_{m,b,n}=T_{m,b,n}\ell_{m,b}.\notag
\end{align}
These costs occur only in the guarantee. The learner does not receive $m$, $b$, $R_{m,b}$, or their minimizing teacher.

\begin{theorem}[Adaptive comparison with all smaller networks]
\label{thm:network-main}
For tanh with $H\ge2$, set $r=H$; for ReLU with $H\ge1$, set $r=L$. Use the root calibration \eqref{eq:root-calibration}. With probability at least $1-\delta$, simultaneously for every learner width $N$ and every $\gamma\ge2$, the actual finite posterior satisfies
\begin{align}
 \int R(z_w)d\Pi_{S,\lambda_S}
 \le\inf_{\substack{1\le m\le N\\b\in\N,\ b\ge1}}
 \bigg\{R_{m,b}+C_r\big[&\mathcal A_{m,b,n}R_{m,b}^{1-1/r}a_n^{2/r}
                  +\mathcal A_{m,b,n}^ra_n^2\big]\bigg\}+n^{-4}.
 \label{eq:network-main}
\end{align}
With one hidden ReLU layer the same assertion holds for every $\gamma\ge1$. The ordinary raw posterior mean obeys the same bound.

For one hidden tanh layer, use the certificate minimization \eqref{eq:shallow-calibration}. With the same probability and simultaneously for every width $N$ and every $\gamma\ge1$,
\begin{align}
 \int R(z_w)d\Pi_{S,\lambda_S}
 \le\inf_{\substack{1\le m\le N\\b\in\N,\ b\ge1}}
 \left\{R_{m,b}+32\left[
 a_n\sqrt{\mathcal A_{m,b,n}R_{m,b}}
       +a_n^2\mathcal A_{m,b,n}\right]\right\}+n^{-4}.
 \label{eq:network-shallow}
\end{align}
Again the ordinary posterior mean satisfies the same bound.
\end{theorem}
\begin{proof}
The Gaussian moment, capacity, and comparison lemmas verify the abstract hypotheses simultaneously over the claimed widths and scales. Take the common capacity event with failure probability $\delta/2$. For each integer pair $(m,b)$, Bernstein's inequality for the fixed population minimizer $g_{m,b}$ gives, except on probability
$\delta/[8(m+1)^2(b+1)^2]$,
\begin{equation}
 \widehat R_S(g_{m,b})\le R_{m,b}
 +(U_{m,b}+1)\sqrt{\frac{2R_{m,b}\ell_{m,b}}n}
 +\frac{2(U_{m,b}+1)^2\ell_{m,b}}{3n}.
 \label{eq:teacher-bernstein}
\end{equation}
Indeed its square loss is at most $(U+1)^2$ and its variance at most $(U+1)^2R_{m,b}$. The sum of the allocated probabilities is less than $\delta/2$. There is therefore one measurable event on which the statistical inequality and every comparison concentration statement hold. No union over the real scale parameter is needed: the statistical event concerns only the common function classes, and the comparison inequalities are deterministic for each posterior.

Fix a pair $(m,b)$ with $m\le N$, write $R=R_{m,b}$, $u=\widehat R_S(g_{m,b})$, and first take $r\ge2$. Set
$D=a_n^2\mathcal A_{m,b,n}^r$.
The explicit implementation budgets imply
\[
 U_{m,b}^2\le E_{m,b}^r,
 \qquad (U_{m,b}+1)^2\ell_{m,b}/n\le4D.
\]
For ReLU, the first inequality follows by raising \eqref{eq:intro-relu-energy} to power $L$ and comparing with \eqref{eq:teacher-amplitudes}; for tanh it follows from $E\ge2Hm^2b^2$ and $E\ge1$. The second uses $a_n^2\ge1/n$ and $\ell_{m,b}\ge1$. Thus \eqref{eq:teacher-bernstein} gives
\begin{equation}
 u\le R+\sqrt{8RD}+\tfrac83D\le(\sqrt R+2\sqrt D)^2.
 \label{eq:bernstein-D}
\end{equation}
Apply \eqref{eq:deep-empirical-oracle} with $A=A_{m,b,n}$ and enlarge its comparison cost from $T_{m,b,n}$ to $\mathcal A_{m,b,n}$. With $p=1-1/r\ge1/2$, its remainder is bounded by
$C_r[D^{1/r}u^p+D]$.
Equation \eqref{eq:bernstein-D} gives
$D^{1/r}u^p\le C_r(R^pD^{1/r}+D)$.
Also $\sqrt{RD}\le R^pD^{1/r}+D$, by splitting into $R\ge D$ and $R<D$. Substitution proves the bound in \eqref{eq:network-main} for this pair. Since every pair is valid on the same event, take the indicated infimum.

For shallow tanh, put $D=a_n^2\mathcal A_{m,b,n}$. The same Bernstein argument gives \eqref{eq:bernstein-D}. The empirical oracle is bounded by
$u+4\sqrt{Du}+20D+n^{-4}$.
Using $u-R\le\sqrt{8RD}+8D/3$ and $\sqrt u\le\sqrt R+2\sqrt D$ makes this at most
$R+32[\sqrt{RD}+D]+n^{-4}$.
Take the infimum over the common countable set of cap pairs. Jensen gives both ordinary-mean statements.
\end{proof}

\begin{remark}[Ordinary biases and noninteger coefficient caps]
Any finite coefficient cap is covered by its integer ceiling. A tanh teacher with arbitrary internal and readout biases compiles into the stated architecture at one extra constant neuron per hidden layer and coefficient cap $2B$. Keep that neuron at value $1/2$: its first bias is $\operatorname{arctanh}(1/2)$ and its later self-edge is $2\operatorname{arctanh}(1/2)=\log3$. Each desired bias $b$ is supplied by an edge of weight $2b$ from the preceding constant neuron. For ReLU, a constant neuron of value one gives the same construction at unchanged cap $B\ge1$. These are comparison constructions only; the learner prior remains the specified dense Gaussian. Thus an ordinarily biased width-$m$ teacher requires $N\ge m+1$.
\end{remark}

\subsection{Rates and the role of approximation error}
At fixed $d,H,\delta$, one has $a_n^2=\widetilde O(1/n)$ and
$\mathcal A_{m,b,n}=\widetilde O(E_{m,b})$, with logarithms in sample count and teacher resources retained by \eqref{eq:simultaneous-cost}. The resulting remainders are
\begin{center}
\begin{tabular}{lll}
\toprule
Activation and depth & General comparison risk $R$ & Exactly fitting teacher \\
\midrule
Tanh, $H\ge2$ & $\widetilde O(E R^{1-1/H}n^{-1/H}+E^H/n)$ & $\widetilde O(E^H/n)$ \\
ReLU, $H\ge1$ & $\widetilde O(E R^{1-1/L}n^{-1/L}+E^L/n)$ & $\widetilde O(E^L/n)$ \\
Tanh, $H=1$ & $\widetilde O(\sqrt{ER/n}+E/n)$ & $\widetilde O(E/n)$ \\
\bottomrule
\end{tabular}
\end{center}
The first two rows have $\gamma\ge2$, with the shallow ReLU extension to all $\gamma\ge1$; the last row permits all $\gamma\ge1$. The calibration is defined at every sample count $n\ge2$ and uses no teacher width, coefficient cap, or risk; a bound larger than the trivial loss scale is simply uninformative.

For a requested excess $\varepsilon$, the first two rows give the sufficient sample form
\begin{equation}
 n=\widetilde O\left(E^r\left[\frac{R^{r-1}}{\varepsilon^r}
                                      +\frac1\varepsilon\right]\right),
 \label{eq:adaptive-sample}
\end{equation}
while shallow tanh gives
$\widetilde O(E[R/\varepsilon^2+1/\varepsilon])$.
The teacher costs and dimensional factors are explicit above; the tilde suppresses only logarithms at fixed dimension and depth.

Let $f_0(x)=\E[Y\mid X=x]$ and $\nu=\E\operatorname{Var}(Y\mid X)$. Then
\[
 R(g)=\nu+\|g-f_0\|_{L^2(P_X)}^2.
\]
Consequently \eqref{eq:network-main} also bounds the learner's squared error relative to $f_0$ by the best teacher approximation error plus the displayed penalty evaluated at $\nu$ plus that approximation error. No exactly realizable network or correct Gaussian noise model is required.

At a fixed coefficient cap, all the displayed learners compete with every teacher width $m\le n^\alpha$ for any fixed $\alpha<1/(2H)$. For tanh with both width and cap $n^\alpha$, the range is $\alpha<1/(4H)$, including the shallow case. For ReLU with both resources $n^\alpha$, the range is $\alpha<1/(4H+2)$. These follow directly from the energy formulas and are attained by the same calibration for every admissible exponent and class; the learner is not told $\alpha$. The stronger conclusion is the entire simultaneous oracle inequality, rather than one selected polynomial-width comparison.

\section{Full-energy capacity for every Lipschitz activation}
\label{sec:all-lipschitz-full-energy}
A scalar covering argument removes the bounded-variation condition on the
activation derivative, at the expense of additional logarithms in sample
size. The activation need only be globally Lipschitz.

For a set $F\subset[-1,1]^n$, write
$\mathfrak R(F)=n^{-1}\mathbb E_\epsilon\sup_{f\in F}\sum_i\epsilon_i f_i$
and let $d_n(f,g)=\|f-g\|_2/\sqrt n$. Let $\mathcal L$ be all scalar
one-Lipschitz functions on $[-1,1]$ which vanish at zero.

\begin{lemma}[An unknown scalar Lipschitz link]
\label{lem:unknown-link}
There is an absolute $C$ such that, for every symmetric nonempty
$F\subset[-1,1]^n$ and $n\ge2$,
\[
 \mathfrak R(\{\psi\circ f:f\in F,\ \psi\in\mathcal L\})
 \le C\log^{3/2}(2n)\mathfrak R(F)+C/\sqrt n.
\]
\end{lemma}
\begin{proof}
We may take closures in the compact cube. Let
$G(F)=n^{-1}\mathbb E_g\sup_{f\in F}\sum_i g_i f_i$, with independent
standard Gaussian coordinates. Conditioning on their magnitudes and
using the Rademacher contraction principle gives
\[
 G(F)\le\mathbb E\max_i|g_i|\,\mathfrak R(F)
       \le C\sqrt{\log(2n)}\,\mathfrak R(F).
\]
Gaussian Sudakov minoration, applied to a maximal separated subset,
therefore gives for $0<a\le2$
\[
 \log\mathcal N(a,F,d_n)
 \le CnG(F)^2/a^2
 \le Cn\log(2n)\mathfrak R(F)^2/a^2.
\]
The scalar class $\mathcal L$ has
$\log\mathcal N(a,\mathcal L,\|\cdot\|_\infty)\le C/a$ for
$0<a\le1$. One can discretize a grid of spacing proportional to $a$,
quantize function values at the same scale, and count only sequences
whose successive quantized increments have bounded integer magnitude;
the initial value is fixed at zero. Representatives may be chosen from
the original class in each nonempty cell.

Composition of these two covers yields
\[
 \log\mathcal N(3a,\mathcal L\circ F,d_n)
 \le Cn\log(2n)\mathfrak R(F)^2/a^2+C/a.
\]
The truncated Dudley entropy bound for Rademacher processes now implies
\[
 \mathfrak R(\mathcal L\circ F)
 \le4\alpha+\frac C{\sqrt n}\int_\alpha^2
          \sqrt{\log\mathcal N(a,\mathcal L\circ F,d_n)}\,da.
\]
Choose $\alpha=n^{-1/2}$. The first entropy term integrates as
$C\sqrt{\log(2n)}\mathfrak R(F)\log(2/\alpha)$, while the second
integrates to $C/\sqrt n$. This proves the claim. Gaussian Sudakov
minoration and the truncated Dudley bound are standard Gaussian and
Rademacher process inequalities~\cite{talagrand}; their metrics and normalizations have
been stated explicitly here.
\end{proof}

Let $\mathcal A_0$ be the closed absolute convex hull of the constant
function one and the $d$ coordinate functions on the input cube.
Recursively define
\[
 \mathcal A_\ell=\overline{\operatorname{aconv}}
   \bigl(\{1\}\cup\{\psi\circ f:f\in\mathcal A_{\ell-1},
                                     \psi\in\mathcal L\}\bigr).
\]
These are symmetric classes with envelope one. Absolute convexification
preserves the Rademacher supremum of a symmetric class. Adding the
constant atom increases its expected supremum by at most $1/\sqrt n$.
Lemma~\ref{lem:unknown-link} and the finite-coordinate bound consequently
show, on every input sample,
\begin{equation}
 \widehat{\mathfrak R}_n(\mathcal A_H)
 \le \frac{C^{H+1}\log^{3H/2}(2n)
                      [1+\sqrt{\log(2d+2)}]}{\sqrt n}.
 \label{eq:all-lipschitz-atom-rad}
\end{equation}
The same bound holds for worst-case empirical complexity.

Now let $\phi$ be globally $A$-Lipschitz, $A\ge1$, and put
$\beta=|\phi(0)|$. At a nonzero preactivation budget $t$, the normalized
profile
\[
 \psi_t(u)=\frac{\phi(tu)-\phi(0)}{At},\qquad -1\le u\le1,
\]
belongs to $\mathcal L$. Assign neuron budgets
\[
 r_j^1=\beta+A\sum_i|W^1_{ji}|,\qquad
 r_j^\ell=\beta+A\sum_i|W^\ell_{ji}|r_i^{\ell-1},\qquad
 P=\sum_j|v_j|r_j^H.
\]
A zero budget gives the zero feature. Otherwise its feature divided by
its budget is a convex combination of the appropriately signed constant atom and a
normalized profile applied to an absolute convex combination from the
preceding class. Induction proves $z_w\in P\mathcal A_H$.

For the usual weighted energy
$Q_N=(d/N)\|W^1\|_F^2+\sum_{\ell\ge2}\|W^\ell\|_F^2+N\|v\|^2$,
Cauchy--Schwarz and $\||W|\|_{\rm op}\le\|W\|_F$ give
\[
 \|r^1\|_2/\sqrt N\le\beta+\sqrt2A\sqrt{Q_N},\qquad
 \|r^\ell\|_2/\sqrt N
    \le\beta+A\sqrt{Q_N}\|r^{\ell-1}\|_2/\sqrt N.
\]
Hence, with $T=2+\beta+2A$ and $L=H+1$,
\[
 P\le T^H(1+Q_N)^{L/2}.
\]
Combining this with \eqref{eq:all-lipschitz-atom-rad} gives
\begin{equation}
 \widehat{\mathfrak R}_n(\{z_w:Q_N(w)\le q\})
 \le\frac{C_{H,A,\beta}\log^{3H/2}(2n)
          [1+\sqrt{\log(2d+2)}]}{\sqrt n}(1+q)^{L/2},
 \qquad |z_w|\le T^H(1+q)^{L/2}.
 \label{eq:all-lipschitz-full-energy-rad}
\end{equation}
The classes include all learner widths. For statistical separability,
use rational-parameter approximations of the actual networks, at each
integer width. Continuity of a Lipschitz activation makes these uniformly
dense on the compact input cube. The enlarged atom classes are only a
capacity bound; their finite-sample closures are compact subsets of the
finite-dimensional sample cube.

For the Gaussian comparison, set $E=d(d+1)B^2+2Hm^2B^2$ and $U=T^H(1+E)^{L/2}$. The perturbation proof in Lemma~\ref{lem:bv-comparison} uses only Lipschitzness and the activation offset, so it applies with this $U$, $B_0=1+\beta+2A(\sqrt E+1)$ and $K_g=2LB_0^H$. Its free-energy cost is
\[
 A_{g,n}=E+1+2(d+H+1)\log[4K_g(U+1)n].
\]
A uniform prior-risk bound is $1+2(4A_0^2)^H$, with $A_0=\max\{1,A,\beta\}$. Use these constants and the displayed capacity bound in the abstract theorem. The countable-cap population argument of Theorem~\ref{thm:network-main} remains valid because the statistical constant includes the envelope factor $T^H$, and $(U+1)^2$ is at most a fixed multiple of $T^{2H}(1+E)^L$.

The full-energy Gaussian posterior argument therefore applies with
$r=L$ to every globally Lipschitz activation, without a condition on
its derivative or its value at zero. The additional
$\log^{3H/2}(2n)$ only changes logarithmic factors in the calibration
and oracle. At deep scales $\gamma\ge2$, and at every scale
$\gamma\ge1$ for one hidden layer, the same teacher comparison and
free-energy root give
\[
 R(\widehat f_n)\le R(g)+\widetilde O_{H,A,\beta,d}
   \left(A_g R(g)^{1-1/L}n^{-1/L}+A_g^L/n\right).
\]
The prior remains one fixed dense independent Gaussian architecture.
The bounded-variation derivative proof remains useful because it avoids
these additional depth-dependent logarithmic powers. Tanh's separate
bounded-activation argument can retain its sharper depth exponent.

\subsection{An optional improvement using generic chaining}
The extra sample-size logarithms in Lemma~\ref{lem:unknown-link} can be
removed by using the two-sided majorizing-measure theorem~\cite{talagrand} instead of
Sudakov minoration followed by an entropy integral. We state the
additional imported theorem explicitly: for a compact index set with
its canonical Gaussian metric, the expected supremum of the centered
Gaussian process is bounded above and below, up to absolute constants,
by Talagrand's $\gamma_2$ functional. We also use the generic-chaining
upper bound for a subgaussian process and the entropy-integral upper
bound on $\gamma_2$.

For a symmetric compact $F\subset[-1,1]^n$ containing zero, retain the
metric $d_n$ and normalized Gaussian complexity $G(F)$ above. The scalar
class $\mathcal L$ is compact in the uniform norm by Arzel\`a--Ascoli.
Its entropy estimate gives
\[
 \gamma_2(\mathcal L,\|\cdot\|_\infty)
 \le C\int_0^2\sqrt{\log\mathcal N(a,\mathcal L,\|\cdot\|_\infty)}\,da
 \le C.
\]
Equip $F\times\mathcal L$ with the sum metric
$d_n(f,f')+\|\psi-\psi'\|_\infty$. Admissible chaining nets for the two
factors, shifted by one level before taking their products, give
\[
 \gamma_2(F\times\mathcal L)
 \le C\{\gamma_2(F,d_n)+\gamma_2(\mathcal L,\|\cdot\|_\infty)\}.
\]
The shift is sufficient because the product of two sets of cardinality
at most $2^{2^{k-1}}$ has cardinality at most $2^{2^k}$.
The map $(f,\psi)\mapsto\psi\circ f$ is one-Lipschitz for this sum
metric. The Gaussian process indexed by its image has canonical metric
$d_n/\sqrt n$. Generic chaining and then the majorizing-measure lower
bound for the Gaussian process indexed by $F$ yield
\[
 G(\mathcal L\circ F)
 \le\frac C{\sqrt n}\{\gamma_2(F,d_n)+1\}
 \le C G(F)+C/\sqrt n.
\]
All processes are centered; symmetry and the zero index make the
expected-supremum conventions unambiguous.

Gaussian width is unchanged by the absolute convex hull of a symmetric
set. Adding the constant atom costs at most $1/\sqrt n$.
The recursion for $\mathcal A_H$ therefore gives
\[
 G(\mathcal A_H)
 \le C^{H+1}[1+\sqrt{\log(2d+2)}]/\sqrt n.
\]
Finally $\mathfrak R(F)\le\sqrt{\pi/2}\,G(F)$: condition on Gaussian
signs and apply Jensen to the independent magnitudes, whose common
mean is $\sqrt{2/\pi}$. Thus the factor
$\log^{3H/2}(2n)$ may be omitted from
\eqref{eq:all-lipschitz-full-energy-rad}, at the cost of another absolute
constant to the power $H$. This optional refinement uses a deeper
standard Gaussian-process theorem; the elementary proof above already
establishes all polynomial-learning and activation-scope claims.

\section{A sharper logarithmic bound from derivative variation}
\label{sec:bv-activations}
A signed-hinge decomposition avoids the additional sample-size logarithms in the preceding general Lipschitz argument when the activation derivative has finite total variation. This gives an alternative proof for the standard activations below, including nonzero activation offsets.

Assume $\phi:\R\to\R$ is globally $K$-Lipschitz and its almost-everywhere derivative has a bounded-variation representative $g$ on $\R$. Write
\[
 a=\phi(0),\qquad V_1=\operatorname{TV}(g),\qquad
 M=\max\{1,K+V_1\},\qquad T=2+|a|+2M.
\]
A Lipschitz function is locally absolutely continuous, so its derivative determines all increments by integration. We may take $|g|\le K$, including its one-sided limits. ReLU is allowed: its derivative has one jump and the corresponding distributional second derivative is a unit point mass. The following safe constants illustrate the scope; numerical sharpness is not needed.
\begin{center}
\begin{tabular}{llll}
\toprule
Activation & $\phi(0)$ & Lipschitz bound $K$ & derivative variation bound $V_1$ \\
\midrule
Tanh & $0$ & $1$ & $2$ \\
ReLU & $0$ & $1$ & $1$ \\
GELU, $x\Phi(x)$ & $0$ & $2$ & $3$ \\
Sigmoid, $(1+e^{-x})^{-1}$ & $1/2$ & $1/4$ & $1/2$ \\
Softplus, $\log(1+e^x)$ & $\log2$ & $1$ & $1$ \\
\bottomrule
\end{tabular}
\end{center}
For GELU, $\phi''(x)=(2-x^2)\varphi(x)$, where $\varphi$ is the standard Gaussian density, so $\int|\phi''|\le\E(2+G^2)=3$; also $|\phi'|=|\Phi+x\varphi|\le2$. The other variation bounds follow from monotonicity on each side of the derivative's single maximum, or from a single jump.

\subsection{Normalized activations are mixtures of centered hinges}
For $t\in[-1,1]$, put
\[
 h_t(u)=\operatorname{ReLU}(u-t)-\operatorname{ReLU}(-t),
 \qquad -1\le u\le1.
\]
Each $h_t$ vanishes at zero, is 1-Lipschitz, and is bounded by one. The identity function is $h_{-1}$ on this interval.

\begin{lemma}[A scale-independent signed mixture]
\label{lem:bv-profile}
For every $s>0$, the function
\[
 u\longmapsto\frac{\phi(su)-\phi(0)}{Ms},\qquad u\in[-1,1],
\]
belongs to the closed absolute convex hull of the centered hinges $h_t$, $t\in[-1,1]$.
\end{lemma}
\begin{proof}
Let $\mu=dg$ be the finite signed distributional derivative measure, and let $\nu_s$ be its restriction to $(-s,s)$ pushed forward by $x\mapsto x/s$. The fundamental theorem for bounded-variation functions gives, on $[-1,1]$,
\begin{equation}
 \frac{\phi(su)-\phi(0)}s
   =g(-s+)u+\int_{(-1,1)}h_t(u)\,d\nu_s(t).
 \label{eq:bv-representation}
\end{equation}
One way to check the identity is to differentiate both sides almost everywhere: the derivative on the right is $g(-s+)+\nu_s((-1,u))=g(su)$ except at a null set, and both sides vanish at zero. Endpoint atoms are absorbed into the one-sided slope or do not affect the interval's interior. Continuity extends equality to the endpoints.

The absolute coefficient mass in \eqref{eq:bv-representation} is at most $K+V_1\le M$. Since $u=h_{-1}(u)$, division by $M$ gives the assertion. The integral is approximated by finite signed mixtures uniformly: $t\mapsto h_t$ is continuous in uniform norm and its signed mixing measure is finite. Zero belongs to the absolute convex hull and supplies any unused mass.
\end{proof}

\subsection{A width-independent capacity estimate}
On a list of $n$ inputs, let $\mathcal A_0$ be the closed absolute convex hull of the $d+1$ coordinate sample vectors of $\bar x=(1,x)$. Recursively let
\[
 \mathcal A_\ell=\overline{\operatorname{aconv}}\left(
 \{\mathbf1\}\cup\{h_t(f):f\in\mathcal A_{\ell-1},\ -1\le t\le1\}\right).
\]
Every class is symmetric, contains zero, and lies in $[-1,1]^n$.

\begin{lemma}[Capacity of the normalized feature classes]
\label{lem:bv-feature-capacity}
For $H\ge0$,
\begin{equation}
 \operatorname{Rad}_n(\mathcal A_H)
 \le\frac{2^H}{\sqrt n}
       \left[\sqrt{2\log(2(d+1))}+4\right].
 \label{eq:bv-feature-capacity}
\end{equation}
\end{lemma}
\begin{proof}
Write $\operatorname{Rad}_n(\mathcal B)=n^{-1}\E_\epsilon\sup_{f\in\mathcal B}\sum_i\epsilon_if_i$. For a symmetric class this also equals the version with an absolute value. The starting class is the convex hull of $2(d+1)$ vectors of Euclidean norm at most $\sqrt n$, so its complexity is at most $\sqrt{2\log(2(d+1))/n}$.

Let $\mathcal A$ be any symmetric class as above and $\mathcal H(\mathcal A)=\{h_t(f):f\in\mathcal A,t\in[-1,1]\}$. Scalar Rademacher contraction for ReLU, with constant one for the ordinary supremum, gives
\[
 \operatorname{Rad}_n(\{\operatorname{ReLU}(f-t):f\in\mathcal A,t\in[-1,1]\})
 \le\operatorname{Rad}_n(\mathcal A)+\E|\textstyle\sum_i\epsilon_i|/n.
\]
The subtracted scalar ranges over $[-1,0]$. Its Rademacher complexity is $\E|\sum_i\epsilon_i|/(2n)$. Splitting the supremum therefore gives
\[
 \operatorname{Rad}_n(\mathcal H(\mathcal A))
 \le\operatorname{Rad}_n(\mathcal A)+\frac3{2\sqrt n}.
\]
This class contains zero. Its positive and negative suprema are nonnegative and have equal expectations, so absolute convexification costs at most a factor two. Adding the constant vector costs at most $1/\sqrt n$. Thus
\[
 \operatorname{Rad}_n(\mathcal A_\ell)
 \le2\operatorname{Rad}_n(\mathcal A_{\ell-1})+4/\sqrt n.
\]
Iterating proves \eqref{eq:bv-feature-capacity}. All classes on a fixed finite input list are compact after closure, so no separability issue enters these finite-dimensional suprema.
\end{proof}

Use the usual full weighted energy
\[
 Q_N=\frac dN\|W^1\|_F^2+\sum_{\ell=2}^H\|W^\ell\|_F^2+N\|v\|_2^2,
 \qquad L=H+1.
\]
\begin{theorem}[A common energy-ball capacity for all BV-derivative activations]
\label{thm:bv-capacity}
For every $q\ge0$, let $\mathcal Z_q$ be the union of raw network outputs over all widths with $Q_N\le q$. Then
\begin{align}
 \sup_{z\in\mathcal Z_q}\|z\|_\infty
 &\le T^H(1+q)^{L/2},\notag\\
 \operatorname{Rad}_n(\mathcal Z_q)
 &\le(2T)^H\left[\sqrt{2\log(2(d+1))}+4\right]
       \frac{(1+q)^{L/2}}{\sqrt n}.
 \label{eq:bv-capacity}
\end{align}
There is no dependence on learner width in either bound.
\end{theorem}
\begin{proof}
Define nonnegative coordinate budgets
\[
 \rho^0=\mathbf1_{d+1},\qquad
 \rho^\ell=|a|\mathbf1_N+M|W^\ell|\rho^{\ell-1}.
\]
An induction shows that hidden neuron $j$'s sample vector lies in $\rho_j^\ell\mathcal A_\ell$. Indeed its preactivation is $s f$ with $s=(|W^\ell|\rho^{\ell-1})_j$ and $f\in\mathcal A_{\ell-1}$. If $s>0$, Lemma~\ref{lem:bv-profile} and the constant component $a$ put its activation in $(|a|+Ms)\mathcal A_\ell$. If $s=0$, its activation is exactly $a$, and the same conclusion holds. A zero budget gives a zero neuron. The output sample vector consequently belongs to $P\mathcal A_H$, where $P=|v|^{\mathsf T}\rho^H$.

Put $b_\ell=\|\rho^\ell\|_2/\sqrt N$. The full energy gives
\[
 b_1\le|a|+M\sqrt{(d+1)/d}\sqrt{Q_N},\qquad
 b_\ell\le|a|+M\sqrt{Q_N}\,b_{\ell-1}.
\]
Since $d\ge1$ and $T\ge|a|+2M$, induction yields
$b_\ell\le T^\ell(1+Q_N)^{\ell/2}$. Therefore
\[
 P\le\sqrt N\|v\|_2 b_H
 \le T^H(1+Q_N)^{L/2}.
\]
The amplitude follows from $\mathcal A_H\subset[-1,1]^n$, and the Rademacher estimate from Lemma~\ref{lem:bv-feature-capacity}. The argument is uniform over every input list and every width. Countable rational parameter approximations with strict energy inequalities, followed by continuity, give a countable determining subclass of the actual network union.
\end{proof}

\subsection{Gaussian teacher neighborhoods without a zero-origin assumption}
At canonical Gaussian scale $s=N^{-\gamma}$, the first, internal, and readout variances are $Ns/d,s,s/N$. Assume $Ns\le1$ and $sD_N\le D_0$, where $D_N=N(d+2)+(H-1)N^2$ and $D_0=d+H+1$. These inequalities hold for every $\gamma\ge2$, and for every $\gamma\ge1$ when $H=1$.

Balanced cloning of a width-$m$, coefficient-$B$ teacher works for every activation, including nonzero offsets: every neuron belongs to a nonempty clone group and has exactly its teacher preactivation. It gives
\[
 E=d(d+1)B^2+2Hm^2B^2
\]
as an upper bound on full weighted energy. By Theorem~\ref{thm:bv-capacity}, take $U=T^H(1+E)^{L/2}$ as an output amplitude bound.

\begin{lemma}[A comparison measure for a general Lipschitz activation]
\label{lem:bv-comparison}
Put
\[
 B_0=1+|a|+2\max\{1,K\}(\sqrt E+1),\qquad
 K_g=2L B_0^H,\qquad
 A_{g,n}=E+1+2D_0\log[4K_g(U+1)n].
\]
At every eligible width and scale, and every $\lambda\ge1/n$,
\begin{equation}
 F_S(\lambda)\le\widehat R_S(g)+\lambda A_{g,n}.
 \label{eq:bv-comparison}
\end{equation}
\end{lemma}
\begin{proof}
Perturb the exact clone by a Gaussian with covariance $\varepsilon^2$ times the prior covariance. Write $b_\ell$ for the normalized $L^2$ size of its random hidden state and $\Delta_\ell$ for its normalized $L^2$ difference from the exact state. Independence of the current Gaussian perturbation, the inequality $Ns\le1$, and $|\phi(t)|\le|a|+K|t|$ give
\[
 b_1\le|a|+\sqrt2K(\sqrt E+\varepsilon),\qquad
 b_\ell\le|a|+K(\sqrt E+\varepsilon)b_{\ell-1}.
\]
For $0<\varepsilon\le1$, these imply $b_\ell\le B_0^\ell$. Lipschitzness also gives
\[
 \Delta_1\le\sqrt2K\varepsilon,\qquad
 \Delta_\ell\le K\sqrt E\Delta_{\ell-1}+K\varepsilon b_{\ell-1}
 \le2\max\{1,K\}\ell\varepsilon B_0^{\ell-1}.
\]
The readout contributes at most $\sqrt E\Delta_H+\varepsilon b_H\le\varepsilon K_g$ to output root mean square error. Only $s\le1$ is used for the readout noise term.

Choose $\varepsilon=[4K_g(U+1)n]^{-1}$. The expected empirical loss is at most $\widehat R_S(g)+1/n$. The exact Gaussian relative-entropy formula gives normalized cost at most
\[
 2s\operatorname{KL}(\rho\Vert\mu)
 \le E+2D_0\log(1/\varepsilon).
\]
Substitute this comparison Gaussian into the free-energy variational identity and use $1/n\le\lambda$. This proves \eqref{eq:bv-comparison}.
\end{proof}

\subsection{The resulting joint activation oracle}
The preceding two results verify the full-energy Gaussian free-energy theorem with
\[
 r=L,\qquad V=T^H,\qquad
 \kappa=(2T)^H\left[\sqrt{2\log(2(d+1))}+4\right],
 \qquad D_0=d+H+1.
\]
A uniform prior-risk bound is also available: with $A_0=\max\{1,K,|a|\}$, conditional Gaussian second moments give
$\E_\mu\widehat R_S\le1+2(4A_0^2)^H$.
Thus the same single free-energy root, with these constants, yields
\begin{equation}
 R(\widehat f_n)\le\inf_g\left\{R(g)+\widetilde O\left(
 A_{g,n}R(g)^{1-1/L}n^{-1/L}+A_{g,n}^L/n\right)\right\}.
 \label{eq:bv-oracle}
\end{equation}
The precise simultaneous interpretation uses the population minimizer of each integer width and coefficient cap and a countable Bernstein union, as in the Gaussian theorem. The prediction is the ordinary mean of raw posterior outputs, and the stronger posterior-average risk obeys the same bound. The event is simultaneous over all eligible widths and every $\gamma\ge2$, or every $\gamma\ge1$ at one hidden layer. The learner knows the activation and architectural depth, but not the comparison cap or risk. No first-layer zero-origin assumption, support restriction, or architecture mixture is used.

At fixed dimension and depth, $A_{g,n}=\widetilde O(E)$, so an exactly fitting teacher gives $\widetilde O(E^L/n)$. Dimension enters polynomially through the explicit comparison and capacity constants. The factors exponential in $H$ are displayed; this argument does not assert polynomial dependence jointly on width and depth. Tanh's sharper depth exponent and ReLU's smaller homogeneous implementation budget remain useful refinements of this common argument.

\section{Extending the argument to every $\gamma\ge1$}
\label{ag-sec:general-lipschitz-raw}

Fix a scalar activation $\phi$ and a constant $A\geq1$ such that
$\operatorname{Lip}(\phi)\leq A$ and $|\phi(0)|\leq A$.
The activation need not be bounded, homogeneous, odd, or zero at the
origin.  In particular the argument covers tanh, ReLU, GELU, sigmoid,
and softplus, with their appropriate fixed Lipschitz constants.

There are $H\geq1$ hidden layers of width $N$, a scalar linear readout,
and a first-layer constant input.  Internal layers have no biases.
Inputs belong to $[-1,1]^d$, labels to $[-1,1]$, and their joint IID
law is arbitrary.  Write $\tau\in(0,1]$ for the actual Gaussian scale:
first-layer, internal, and readout coordinate variances are
$\tau^2/d$, $\tau^2/N$, and $\tau^2/N^2$.
Put $\omega=N/\tau^2$.  The likelihood variance is $n\lambda/\omega$,
so the actual posterior is
\[
 d\Pi=Z^{-1}\exp[-\omega\widehat R_S(z_w)/(2\lambda)]\,d\mu.
\]
This is the common prior and posterior in the notation $s=\tau^2/N=1/\omega$. Taking $\tau^2=N^{1-\gamma}$ gives every canonical prior exponent $\gamma\ge1$. The additional $\tau$ notation makes the width-uniform estimates explicit.

\subsection{A teacher cylinder without bounded activations}

Let $g$ be a width-$m$ teacher, with all coefficients at most $B\geq1$
in absolute value.  Define the weighted Gaussian energy
\[
 Q_N(w)=\frac dN\|W^1\|_F^2+
             \sum_{\ell=2}^H\|W^\ell\|_F^2+N\|v\|_2^2.
\]
The Gaussian density has exponent $-\omega Q_N(w)/2$.
To represent the teacher at any $N\geq m$, partition each hidden
layer into clone groups of sizes
$k_j\in\{\lfloor N/m\rfloor,\lceil N/m\rceil\}$.
Copy each teacher first-layer row to its group. For an internal
teacher edge from neuron $i$ to neuron $j$, assign its weight divided
by $k_i$ to every edge between the corresponding clone groups.
Divide each teacher readout weight by its group size in the same way.
The sums entering every clone equal the teacher sums, so this network
computes exactly $g$, for any activation.

The first-layer contribution is at most $d(d+1)B^2$.
Since $k_j/k_i\leq2$, each internal-layer contribution is at most
$2m^2B^2$. Since $N/k_i\leq2m$, the readout contribution is at most
$2m^2B^2$. Thus the clone $w_*$ satisfies
\[
 Q_N(w_*)\leq E:=d(d+1)B^2+2Hm^2B^2.
\]
Every learner neuron belongs to a nonempty clone group, so this
construction does not require $\phi(0)=0$.
Define
\begin{gather*}
 b_0=2\sqrt{2[2\log9+\log(4(H+1))]},\qquad
 T_g=A[3+2mB+B(d+1)+b_0],\\
 P_g=1+(\sqrt2mB+1)A(d+1+H\sqrt m)T_g^{H-1}+T_g^H,\\
 D_g=d+2+(H-1)m,\qquad
 \eta_\rho=\min\{1/2,\rho/(2P_g)\},\qquad c_0=3/8,\\
 \mathcal A_g(\rho)=E/2+D_g\log(1/(c_0\eta_\rho))+\log2.
\end{gather*}
Let $U\geq\|g\|_\infty$ be any specified output bound;
one may always take $U=mBT_g^H$, and for tanh one may take $U=mB$.

\begin{lemma}[A finite-width teacher neighborhood]
\label{ag-lem:general-activation-cylinder}
For every $N\geq m$, $0<\tau\leq1$, and $0<\rho\leq1$,
\[
 -\log\mu\{\|z_w-g\|_\infty<\rho\}
 \leq\omega E/2+ND_g\log(1/(c_0\eta_\rho))+\log2
 \leq\omega\mathcal A_g(\rho).
\]
\end{lemma}
\begin{proof}
Let $w_*$ be the balanced clone.  At each hidden layer let $P$ project
onto the $m$-dimensional subspace of vectors constant on clone groups.
In standardized independent Gaussian coordinates, require absolute
value at most $\eta_\rho$ for all first-layer and readout coordinates,
and for every internal row's projections onto an orthonormal basis of
this group subspace.  There are $ND_g$ independent scalar constraints,
whose joint probability is at least $(c_0\eta_\rho)^{ND_g}$.

Conditioning on these intervals preserves independent centered
unit-subgaussian coordinates in the corresponding orthonormal bases.
Two sphere nets, each of size at most $9^N$, give for an internal
standardized matrix $Z$
\[
 \Pr\{\|Z\|_{\rm op}>b_0\sqrt N\mid\text{cylinder}\}
 \leq2\exp[(2\log9-b_0^2/8)N].
\]
A union bound over internal layers leaves a centrally symmetric event
of probability at least one half the cylinder probability.

Translate this event by $w_*$.  The internal perturbation has operator
norm at most $b_0$, while its restriction to the teacher's group
subspace has operator norm at most $\eta_\rho\sqrt m$.
The clone internal operator norm is at most $\sqrt2mB$.
Teacher hidden coordinates are at most $T_g^\ell$ at layer $\ell$,
because $|\phi(s)|\leq A(1+|s|)$.
For normalized hidden errors $e_\ell=\|h^\ell-h_*^\ell\|_2/\sqrt N$,
Lipschitz propagation gives
\[
 e_1\leq A\eta_\rho(d+1),\qquad
 e_\ell\leq A(\sqrt2mB+b_0)e_{\ell-1}
                +A\eta_\rho\sqrt m\,T_g^{\ell-1}.
\]
Consequently
$e_H\leq\eta_\rho A(d+1+H\sqrt m)T_g^{H-1}$.
The clone readout has $\sqrt N\|v_*\|_2\leq\sqrt2mB$, and its
perturbation has $\sqrt N\|v-v_*\|_2\leq\eta_\rho$.
Thus the output error is at most
$(\sqrt2mB+1)e_H+\eta_\rho T_g^H<\rho$.

For a centered Gaussian measure and a centrally symmetric event,
translation by a vector of squared standardized length $r$ decreases
its probability by at most the factor $e^{-r/2}$; this follows by
pairing opposite points and applying Jensen to the Gaussian density
ratio.  Here $r\leq\omega E$.  This proves the first bound; $N\leq\omega$
and $1\leq\omega$ prove the second.
\end{proof}

\subsection{Width-uniform prior moments and norm tails}

Define
\[
 K(w)=\max\{128,\|W^1\|_F^2/N,
          \|W^2\|_{\rm op}^2,\ldots,\|W^H\|_{\rm op}^2,
          N\|v\|_2^2\}.
\]
Gaussian chi-square tails for the first layer and readout, and two
sphere nets for internal layers, imply for $u\geq128$
\begin{equation}
 \mu\{K>u\}\leq(H+1)e^{-\omega u/16}.
 \label{ag-eq:general-activation-norm-tail}
\end{equation}
For example the internal estimate is
$2\exp[2N\log9-\omega u/8]\leq e^{-\omega u/16}$.
The first-layer term in $K$ is $(\tau^2/(dN))\chi^2_{N(d+1)}$;
the readout term is $(\tau^2/N)\chi^2_N$, so the same weaker bound
applies to them.
Tail integration yields
\begin{equation}
 \log\int e^{\omega K/32}d\mu\leq\omega B_H,
 \qquad B_H=4+\log(H+2).
 \label{ag-eq:general-activation-norm-mgf}
\end{equation}

There are also width-uniform raw-output moments.  Let
$S_\ell=\|h^\ell(x)\|_2^2/N$, and use $S_0=\|\bar x\|_2^2/d\leq2$
for the first conditional Gaussian field.  Conditional on previous
layers, each new field has variance $\tau^2S_{\ell-1}$.
The bounds $|\phi(s)|\leq A(1+|s|)$ and Gaussian second and fourth
moments give
\[
 \mathbb E S_\ell\leq2(4A^2)^\ell,
 \qquad \mathbb E S_\ell^2\leq4(32A^4)^\ell.
\]
For the second assertion, Jensen applied to the average of squared
coordinates and $|\phi(s)|^4\leq8A^4(1+|s|^4)$ suffice; these bounds
also apply to the first layer with the stated $S_0$.
The centered Gaussian readout therefore implies, uniformly in $x$,
\[
 \mathbb E_\mu z_w(x)^2\leq 2(4A^2)^H/N,
 \qquad \mathbb E_\mu z_w(x)^4\leq12(32A^4)^H/N^2.
\]
Put
\[
 R_0=1+2(4A^2)^H,\qquad V_0=8+96(32A^4)^H.
\]
Tonelli and Jensen give, for every sample and every data law under
consideration,
\begin{equation}
 \int\widehat R_S(z_w)d\mu\leq R_0,\qquad
 \int R(z_w)d\mu\leq R_0,\qquad
 \int R(z_w)^2d\mu\leq V_0,
 \qquad Z\geq e^{-\omega R_0/(2\lambda)}.
 \label{ag-eq:general-activation-prior-moments}
\end{equation}

\subsection{A common function cover and raw statistical transfer}

Write $L=H+1$.  On $K(w)\leq k$, Lipschitz propagation gives
\[
 \operatorname{Lip}(z_w)\leq A^H k^{L/2},\qquad
 \|z_w\|_\infty\leq C_{d,H,A}k^{L/2}.
\]
For the amplitude bound, normalized hidden norms satisfy
$a_\ell\leq A(1+\sqrt k\,a_{\ell-1})$ starting from
$a_0\leq\sqrt{d+1}$, and the readout multiplies by at most $\sqrt k$.
Thus this bound includes nonzero $\phi(0)$.
An input grid of mesh proportional to $a/k^{L/2}$, followed by rounding
values to multiples of $a$, provides a deterministic finite cover
$\mathcal C_{k,a}$ in uniform norm, common to all widths and scales,
with centers bounded by $C_{d,H,A}k^{L/2}+1$ and
\begin{equation}
 \log|\mathcal C_{k,a}|
 \leq C_{d,H,A} k^{dL/2}a^{-d}
                  \log\frac{C_{d,H,A}k}{a}
 \qquad(k\geq128,\ 0<a<1).
 \label{ag-eq:general-activation-cover}
\end{equation}
Nearest-grid-point functions are sufficient; the cover centers need
not themselves be networks.


\section{A statistical bound for finite covers}
\label{ag-sec:finite-cover-statistics}

The only uniform statistical estimate needed below is an elementary
finite-class consequence of Bernstein's inequality. The empirical
risk, rather than a worst-case risk bound, appears in its leading term.

\begin{lemma}[Relative deviation for a finite prediction class]
\label{ag-lem:finite-cover-bernstein}
Let $\mathcal F$ be a finite class of measurable real functions, all
bounded in absolute value by $M\geq1$, and let $|Y|\leq1$.
Put $b=(M+1)^2$ and $t=\log(|\mathcal F|/\delta)$.
For $n$ IID observations, with probability at least $1-\delta$,
simultaneously for all $h\in\mathcal F$,
\begin{equation}
 R(h)\leq\widehat R_S(h)
       +3\sqrt{\frac{b\widehat R_S(h)t}{n}}
       +\frac{4bt}{n}.
 \label{ag-eq:finite-cover-bernstein}
\end{equation}
\end{lemma}
\begin{proof}
For one $h$, the squared loss $\ell=(h(X)-Y)^2$ lies in $[0,b]$,
and $\operatorname{Var}(\ell)\leq\mathbb E\ell^2\leq bR(h)$.
The one-sided Bernstein bound and a union bound give
\[
 R(h)-\widehat R_S(h)
 \leq\sqrt{2bR(h)t/n}+2bt/(3n)
\]
simultaneously over the class. Set
$x=(R(h)-\widehat R_S(h))_+$ and $a=2bt/n$.
Then
\[
 x\leq\sqrt{a(\widehat R_S(h)+x)}+a/3
 \leq\sqrt{a\widehat R_S(h)}+x/2+a/2+a/3.
\]
Rearranging gives
$x\leq2\sqrt{2b\widehat R_S(h)t/n}+10bt/(3n)$,
which implies the displayed weaker constants.
\end{proof}

If a network output is uniformly within $a$ of a cover center, then
both population and empirical square-root risks differ by at most $a$,
by Minkowski's inequality. In the proofs below the empirical risk is
bounded by a constant. The cover estimate first bounds the center's
population risk by a constant, after which this square-root comparison
transfers its deviation to the network with a cost proportional to $a$.

\section{Free-energy calibration and one statistical event}
\label{ag-sec:adaptive-core}
For one width and prior scale define the empirical free energy
\[
 F_S(\lambda)=-\frac{2\lambda}{\omega}\log
 \int e^{-\omega\widehat R_S(z_w)/(2\lambda)}d\mu(w).
\]
The Gibbs variational identity gives
\begin{equation}
 F_S(\lambda)=\inf_\rho\left\{\int\widehat R_S(z_w)d\rho
                    +\frac{2\lambda}{\omega}\KL(\rho\Vert\mu)\right\}.
 \label{ag-eq:variational}
\end{equation}
In particular $0\le F_S(\lambda)\le R_0$. We suppress $S$ where it causes no ambiguity. The Gaussian likelihood's parameter-independent normalizing factor is omitted in this definition. The calibration below is therefore not maximum marginal-likelihood selection of a Gaussian noise variance.

Put $\rho_d=(d+2)/2$. There are two valid choices of exponents:
\begin{center}
\begin{tabular}{lll}
\toprule
Activation & $p$ & $q$ \\
\midrule
Any fixed globally Lipschitz $\phi$ & $L\rho_d$ & $(L-1)\rho_d$ \\
Tanh refinement & $\max\{d+2,1+dL/2\}$ & $p-\rho_d$ \\
\bottomrule
\end{tabular}
\end{center}
In both cases let $r=p/\rho_d\ge2$; thus $q/p=1-1/r$.
Define
\begin{align}
 C_{\rm tail}&=R_0+\sqrt{V_0(H+1)},&
 b_n&=\max\{16,4\log(C_{\rm tail}n^3)\},\notag\\
 U_0&=2R_0+1,&
 \theta_n&=C_{d,H,A}\frac{\log^6(2n/\delta)}n.
 \label{ag-eq:calibration-constants}
\end{align}
Here and below $C_{d,H,A}$ is a sufficiently large fixed constant depending only on the displayed parameters. Its choice is fixed before sampling. We initially suppose
\begin{equation}
 \theta_n b_n^p U_0^q\le1.
 \label{ag-eq:admissibility}
\end{equation}
For fixed $d,H,A,\delta$, this holds at every sufficiently large sample count. A final remark supplies a bounded-risk fallback when it fails.

\begin{lemma}[One measurable likelihood calibration]
\label{ag-lem:root}
For every sample, width, and scale, there is a unique positive solution of
\begin{equation}
 \lambda^p=\theta_n[2F_S(\lambda)+\lambda b_n]^q.
 \label{ag-eq:root}
\end{equation}
This solution is measurable in the sample. Under \eqref{ag-eq:admissibility}, putting $u=2F_S(\lambda)+\lambda b_n$, it satisfies
\begin{equation}
 1/n\le\lambda\le1/b_n,\qquad 16\lambda\le u\le U_0.
 \label{ag-eq:root-range}
\end{equation}
\end{lemma}
\begin{proof}
The variational representation shows that $F$ is finite, concave, and continuous on $(0,\infty)$. Its Laplace representation shows that $F(\lambda)/\lambda$ is nonincreasing. Consequently
\[
 \Psi(\lambda)=\frac{\lambda^p}{[2F(\lambda)+\lambda b_n]^q}
 =\frac{\lambda^{\rho_d}}{[2F(\lambda)/\lambda+b_n]^q}
\]
is continuous and strictly increasing. Its limit at zero is zero, because it is at most $\lambda^{\rho_d}/b_n^q$; its limit at infinity is infinity, because $F\le R_0$. The positive root is unique. For each fixed $t>0$, the event $\{\lambda<t\}$ is $\{\Psi_S(t)>\theta_n\}$, proving measurability.

At $\lambda=1/b_n$, the bracket in the first denominator is at most $U_0$, so \eqref{ag-eq:admissibility} gives $\Psi(1/b_n)\ge\theta_n$. Thus the root is at most $1/b_n$. Its equation also implies $\lambda^{\rho_d}\ge\theta_n b_n^q\ge1/n$. Since $\rho_d\ge1$, this implies $\lambda\ge1/n$. The bounds on $u$ now follow.
\end{proof}

\begin{lemma}[A weighted joint tail]
\label{ag-lem:joint-tail}
For the root in Lemma~\ref{ag-lem:root}, define
\[
 J_\lambda(w)=\widehat R_S(z_w)+\lambda K(w)/16.
\]
Then, deterministically for every sample,
\begin{equation}
 \int R(z_w)\one_{\{J_\lambda>u\}}d\Pi_{S,\lambda}(w)\le n^{-3}.
 \label{ag-eq:joint-tail}
\end{equation}
\end{lemma}
\begin{proof}
If $J_\lambda>u$, then either $\widehat R_S>u/2$ or $K>8u/\lambda$. The latter threshold is at least128 by \eqref{ag-eq:root-range}. On the first event, the likelihood is at most $e^{-\omega u/(4\lambda)}$, so its posterior raw-risk integral is at most
\[
 R_0\exp\left[\frac{\omega F(\lambda)}{2\lambda}
                      -\frac{\omega u}{4\lambda}\right].
\]
On the second event, bound the likelihood by one and apply Cauchy--Schwarz to the prior. The prior fourth-moment estimate and norm tail give an upper bound
\[
 \sqrt{V_0(H+1)}\exp\left[\frac{\omega F(\lambda)}{2\lambda}
                      -\frac{\omega u}{4\lambda}\right].
\]
Since $u=2F(\lambda)+\lambda b_n$ and $\omega\ge1$, their sum is at most $C_{\rm tail}e^{-b_n/4}\le n^{-3}$. This argument neither bounds the full parameter dimension nor replaces the posterior by a norm-restricted learner.
\end{proof}

\begin{lemma}[A common empirical-loss and norm certificate]
\label{ag-lem:common-statistical}
For a sufficiently large choice of $C_{d,H,A}$ in \eqref{ag-eq:calibration-constants}, there is an event of probability at least $1-\delta/2$ on which the following assertion holds simultaneously for all widths, all prior scales, and every positive $\lambda,u$ satisfying
\[
 1/n\le\lambda\le1/b_n,\qquad16\lambda\le u\le U_0,
 \qquad\lambda^p\ge\theta_n u^q.
\]
Every parameter vector with $J_\lambda(w)\le u$ satisfies
\begin{equation}
 R(z_w)\le\widehat R_S(z_w)+\lambda K(w)/16.
 \label{ag-eq:local-statistical}
\end{equation}
\end{lemma}
\begin{proof}
Construct dyadic norm levels $k_j=128\cdot2^j$ through a level at least $32U_0n$, and dyadic uniform accuracies through a level below $c_1/n$, beginning at $1/2$. Constants $c_1>0$ below depend only on $d,H,A$. There are $O_{d,H,A}(\log^2(2n))$ pairs. For each pair use the deterministic cover from \eqref{ag-eq:general-activation-cover}, and apply Lemma~\ref{ag-lem:finite-cover-bernstein} with its share of failure probability $\delta/2$. This defines one event from fixed finite classes before the sample is observed.

Fix now any admissible $\lambda,u$ and any $w$ with $J_\lambda(w)\le u$. Choose a norm level
\[
 K(w)\le k\le2K(w)\le32u/\lambda.
\]
Choose a grid accuracy $a$ within a factor two of
\[
 a_* = c\lambda k/\sqrt u,
\]
where $c>0$ is sufficiently small. Since $k\ge128$, $\lambda\ge1/n$, and $u\le U_0$, this accuracy exceeds $c_1/n$ after choosing $c_1$ small enough. Since $\lambda k\le32u$, one also has $a_*\le32c\sqrt{U_0}<1/2$. Thus the chosen pair occurs in the deterministic list.

Let $h$ be a cover center within $a$ of $z_w$. Its empirical risk obeys
$\sqrt{\widehat R_S(h)}\le\sqrt u+a\le C\sqrt u$.
Its square loss is bounded by $C k^\beta$, with $\beta=L$ for a general activation and $\beta=1$ for tanh. In the latter case, $|z_w|\le\sqrt K$ because tanh hidden outputs are bounded by one; the same covers and Lipschitz estimates still apply. Write
\[
 a_0=\beta+dL/2,
 \qquad T_n=\log(2n/\delta).
\]
The cover bound and the finite allocation of failure probabilities give
\[
 \frac{k^\beta}{n}\log\frac{|\mathcal C_{k,a}|}{\eta_{k,a}}
 \le C\frac{T_n^6}{n}\,
        u^{d/2}\lambda^{-d}k^{a_0-d}.
\]
The deliberately generous power six also absorbs all fixed constants and grid-count logarithms. Hence the square-root Bernstein term is at most $\varepsilon\lambda k$, for any fixed sufficiently small $\varepsilon>0$, provided
\begin{equation}
 n\ge C_{d,H,A}T_n^6
      u^{\rho_d}\lambda^{-2\rho_d}k^{a_0-2\rho_d}.
 \label{ag-eq:shell-condition}
\end{equation}
The linear Bernstein term is then at most $C\varepsilon^2\lambda k$, since $\lambda k\le32u$. It also follows that $R(h)\le C u$. Minkowski for population and empirical root risks therefore transfers the estimate from $h$ to $z_w$ with extra cost at most
\[
 C(a\sqrt u+a^2)\le Cc\lambda k.
\]
Choose $c,\varepsilon$ small enough that the total cost is at most $\lambda k/32\le\lambda K(w)/16$.

It remains to verify \eqref{ag-eq:shell-condition}. In the general case $a_0=L\rho_d=p\ge2\rho_d$. Using $k\le32u/\lambda$, its right side is bounded by
$C_{d,H,A}T_n^6 u^{p-\rho_d}\lambda^{-p}$.
For tanh, if $a_0\ge2\rho_d$ the same argument applies with $p=a_0$. If $a_0<2\rho_d$, use the constant lower bound $k\ge128$ instead, obtaining $C_{d,H,A}T_n^6u^{\rho_d}\lambda^{-2\rho_d}$; here $p=2\rho_d$ and $q=\rho_d$. In every case the required inequality is exactly implied by $\lambda^p\ge\theta_nu^q$, after fixing the constant in $\theta_n$ sufficiently large.
\end{proof}

\begin{theorem}[Free-energy certificate at every prior scale]
\label{ag-thm:certificate}
On the event in Lemma~\ref{ag-lem:common-statistical}, simultaneously for every width $N$ and every $0<\tau\le1$, the root-calibrated actual finite posterior satisfies
\begin{equation}
 \int R(z_w)d\Pi_{S,\lambda}(w)
 \le F_S(\lambda)+2B_H\lambda+n^{-3}.
 \label{ag-eq:certificate}
\end{equation}
Its ordinary raw posterior mean has the same risk bound.
\end{theorem}
\begin{proof}
Integrate \eqref{ag-eq:local-statistical} over $J_\lambda\le u$, extend the nonnegative empirical-loss and norm terms to the whole parameter space, and use \eqref{ag-eq:joint-tail}. The prior exponential moment \eqref{ag-eq:general-activation-norm-mgf} and the entropy inequality give
\[
 \frac{\lambda}{16}\int Kd\Pi
 \le\frac{2\lambda}{\omega}\KL(\Pi\Vert\mu)+2B_H\lambda.
\]
The variational identity at its posterior optimizer proves \eqref{ag-eq:certificate}. The raw-risk integral is finite, and Jensen proves the ordinary-mean claim. Every estimate other than the common statistical event is deterministic, so the event is simultaneous over the uncountable scale parameter as stated.
\end{proof}

\section{Simultaneous comparison with unknown teacher classes}
\label{ag-sec:population-oracle}
For positive integers $m,b$, let $\mathcal G_{m,b}$ be the depth-$L$ networks of width at most $m$ and coefficient cap $b$ in the architecture above. They can be represented as a compact finite-dimensional parameter set by zero padding unused neurons and setting their outgoing weights to zero. This works even when $\phi(0)\ne0$. Their outputs are uniformly bounded, so population square loss is continuous on that set. A deterministic population minimizer $g_{m,b}$ exists. Write
\[
 R_{m,b}=\min_{g\in\mathcal G_{m,b}}R(g)\le1,
\]
since the zero readout gives the zero function.

Use the definitions $E,T_g,P_g,D_g,\mathcal A_g$ from Section~\ref{ag-sec:general-lipschitz-raw} with $B=b$, and put $U_{m,b}=mbT_g^H$; for tanh use the sharper $U_{m,b}=mb$. Define
\begin{align}
 A_{m,b,n}&=1+2\mathcal A_g\left(\frac1{4n(U_{m,b}+1)}\right),\notag\\
 \ell_{m,b}&=1+\log\frac{32(m+1)^2(b+1)^2}{\delta},\notag\\
 \mathcal B_{m,b,n}&=(A_{m,b,n}+b_n+2B_H+1)\ell_{m,b}.
 \label{ag-eq:comparison-cost}
\end{align}
These are fully specified width-independent comparison costs. They are used only in the bound, not in the likelihood calibration.

\begin{lemma}[A comparison valid throughout the selected range]
\label{ag-lem:teacher-comparison}
For every sample, every $N\ge m$, every $0<\tau\le1$, and every $\lambda\ge1/n$,
\begin{equation}
 F_S(\lambda)\le\widehat R_S(g_{m,b})+\lambda A_{m,b,n}.
 \label{ag-eq:teacher-comparison}
\end{equation}
\end{lemma}
\begin{proof}
Condition the prior on the uniform output neighborhood of radius
$\varepsilon=1/[4n(U_{m,b}+1)]$ from Lemma~\ref{ag-lem:general-activation-cylinder}. Throughout this neighborhood,
\[
 \widehat R_S(z_w)\le\widehat R_S(g_{m,b})
                      +2(U_{m,b}+1)\varepsilon+\varepsilon^2
 \le\widehat R_S(g_{m,b})+1/n.
\]
The conditional prior has relative entropy at most $\omega\mathcal A_g(\varepsilon)$. Substituting it into \eqref{ag-eq:variational}, and using $1/n\le\lambda$, proves the claim. A uniform output tube is only a comparison construction; the actual posterior is unrestricted.
\end{proof}

\begin{theorem}[An adaptive oracle throughout $\gamma\ge1$]
\label{ag-thm:main}
Suppose \eqref{ag-eq:admissibility}, and select the likelihood by \eqref{ag-eq:root}. With probability at least $1-\delta$, simultaneously for every learner width $N$ and every $0<\tau\le1$, the actual finite posterior satisfies
\begin{align}
 \int R(z_w)d\Pi_{S,\lambda}(w)
 \le\inf_{\substack{1\le m\le N\\b\in\N,\ b\ge1}}
 \bigg\{R_{m,b}+C_{d,H,A}\big[&
 \mathcal B_{m,b,n} R_{m,b}^{1-1/r}\theta_n^{1/p}
 \notag\\[-2pt]
 &+\mathcal B_{m,b,n}^r\theta_n^{1/\rho_d}\big]\bigg\}+n^{-3}.
 \label{ag-eq:population-oracle}
\end{align}
The same inequality holds for the ordinary raw posterior mean. In particular it holds simultaneously for every canonical exponent $\gamma\ge1$, by $\tau^2=N^{1-\gamma}$.
\end{theorem}
\begin{proof}
Use the statistical event of Lemma~\ref{ag-lem:common-statistical}. For each pair $(m,b)$, apply Bernstein's inequality to the fixed population minimizer $g_{m,b}$, with failure probability $\delta/[8(m+1)^2(b+1)^2]$. Its loss is at most $(U_{m,b}+1)^2$ and its variance at most $(U_{m,b}+1)^2R_{m,b}$. The sum of these failure probabilities is less than $\delta/2$. Thus on one event of probability at least $1-\delta$, all teacher concentrations and the common statistical inequality hold. No union over learner scales or over data-selected roots is required.

Fix one cap pair and a learner with $N\ge m$, and abbreviate $R=R_{m,b}$, $u_g=\widehat R_S(g_{m,b})$, $A_g=A_{m,b,n}$, and $\mathcal B=\mathcal B_{m,b,n}$. The root equation and Lemma~\ref{ag-lem:teacher-comparison} imply
\[
 \lambda^p\le\theta_n[2u_g+(2A_g+b_n)\lambda]^q.
\]
Split according to which summand in the bracket is larger. Taking the $p$th root in the first case and cancelling $\lambda^q$ in the second gives
\begin{equation}
 \lambda\le C_{d,H,A}\left[
 \theta_n^{1/p}u_g^{q/p}
 +\theta_n^{1/\rho_d}(2A_g+b_n)^{q/\rho_d}\right].
 \label{ag-eq:root-comparison}
\end{equation}
Combine this with \eqref{ag-eq:certificate} and the teacher free-energy bound. Since $r=p/\rho_d$ and $q=p-\rho_d$, writing
\[
 D=\mathcal B^r\theta_n^{1/\rho_d}
\]
yields the empirical oracle
\begin{equation}
 \int R(z_w)d\Pi\le u_g+C_{d,H,A}
       [D^{1/r}u_g^{1-1/r}+D]+n^{-3}.
 \label{ag-eq:empirical-oracle}
\end{equation}

The explicit teacher amplitudes obey
$(U_{m,b}+1)^2\le C_{d,H,A}(1+E)^r$.
For general activations, $T_g\le C_{d,H,A}mb$ and $E\ge m^2b^2$, giving this bound with $r=L$. For tanh, $U_{m,b}=mb$ and $r\ge2$ suffice. Since $\theta_n\le1$, $\rho_d\ge1$, $A_g\ge1+E$, and $\ell_{m,b}\ge1$, the quantity $D$ dominates
$C^{-1}_{d,H,A}(U_{m,b}+1)^2\ell_{m,b}/n$.
Bernstein therefore gives
\begin{equation}
 u_g\le R+C_{d,H,A}\sqrt{RD}+C_{d,H,A}D.
 \label{ag-eq:teacher-risk-concentration}
\end{equation}
For $r\ge2$,
\[
 \sqrt{RD}\le R^{1-1/r}D^{1/r}+D,
\]
as follows by splitting into $R\ge D$ and $R<D$. Also \eqref{ag-eq:teacher-risk-concentration} implies $u_g\le C_{d,H,A}(R+D)$, hence
\[
 D^{1/r}u_g^{1-1/r}\le C_{d,H,A}
                  [R^{1-1/r}D^{1/r}+D].
\]
Keep the leading $R$ in \eqref{ag-eq:teacher-risk-concentration} with coefficient one and apply these two inequalities to its remainder and to \eqref{ag-eq:empirical-oracle}. This proves \eqref{ag-eq:population-oracle} for the chosen pair. Every pair is valid on the same event, so take the stated infimum. Jensen proves the ordinary-mean claim.
\end{proof}

\begin{remark}[A definition at every sample count]
If \eqref{ag-eq:admissibility} fails, select $\lambda=\omega R_0$ instead. Jensen gives $Z\ge e^{-1/2}$, so the posterior Gibbs risk is at most $e^{1/2}R_0$, deterministically. Failure of \eqref{ag-eq:admissibility} implies
\[
 \mathcal B_{m,b,n}^r\theta_n^{1/\rho_d}
 \ge b_n^r\theta_n^{1/\rho_d}
 >U_0^{-q/\rho_d}.
\]
Thus, after increasing only $C_{d,H,A}$, \eqref{ag-eq:population-oracle} also holds in this case. This supplies one specified learner for every $n\ge2$. For all sufficiently large $n$, it is the free-energy-calibrated learner of the main theorem.
\end{remark}

\begin{remark}[Ordinary internal biases]
If $\phi$ is not identically zero, choose any fixed $a_0$ with $c=\phi(a_0)\ne0$. One constant neuron per layer can be kept at value $c$, using first bias $a_0$ and later self-edge $a_0/c$. An ordinary internal bias $b$ is then supplied by an edge $b/c$, and the same construction supplies a readout bias. Thus ordinary biased teachers compile at one extra neuron per layer and a coefficient factor depending only on the fixed activation. This changes comparison resources by fixed factors and requires learner width at least $m+1$. If $\phi$ is identically zero, the model as specified has only the zero output, so the theorem is trivial for its own comparison classes.
\end{remark}

\subsection{Sample exponents}
At fixed $d,H,A,\delta$, the costs in \eqref{ag-eq:comparison-cost} obey $\mathcal B_{m,b,n}=\widetilde O(E)$, with logarithmic dependence on the displayed resources. Since $p=r(d+2)/2$, the remainder in \eqref{ag-eq:population-oracle} is
\[
 \widetilde O\left(E R^{1-1/r}n^{-2/[r(d+2)]}
                       +E^r n^{-2/(d+2)}\right).
\]
For a requested excess $\varepsilon>0$, the sufficient sample form is
\[
 n=\widetilde O\left(
 E^{r(d+2)/2}
 \left[\frac{R^{(r-1)(d+2)/2}}{\varepsilon^{r(d+2)/2}}
              +\frac1{\varepsilon^{(d+2)/2}}\right]\right).
\]
The coefficient on comparison risk remains one. If $f_0(x)=\E[Y\mid X=x]$ and $\nu=\E\operatorname{Var}(Y\mid X)$, then $R(g)=\nu+\|g-f_0\|_{L^2(P_X)}^2$. Subtracting $\nu$ turns the theorem into an approximation bound for the conditional mean; the penalty is evaluated at noise plus approximation error. The selected learner is told neither quantity.

\section{A sharper ReLU comparison cost at every scale}
\label{ag-sec:relu-selective}
Positive homogeneity and a sparse comparison mean improve the teacher budget without restricting the learner prior. Only this section assumes ReLU.

\begin{lemma}[A balanced sparse teacher]
\label{ag-lem:relu-balanced}
A width-$m$, coefficient-$B$ teacher has an exact implementation at every $N\ge m$, with at most $m$ active coordinates per hidden layer, and total energy
\[
 E_{\rm R}=L[d(d+1)]^{1/L}m^{2H/L}B^2.
\]
Each layer's weighted energy is at most $e=E_{\rm R}/L$. Its output amplitude is at most $U_{\rm R}=(d+1)m^HB^L$.
\end{lemma}
\begin{proof}
Place the teacher into the first $m$ coordinates of every hidden layer. Multiply the first-layer mean matrix by $\sqrt N$ and the readout by $1/\sqrt N$, preserving the output by ReLU homogeneity. The weighted layer energies become $d\|A_1\|_F^2,\|A_2\|_F^2,\ldots,\|A_H\|_F^2,\|c\|_2^2$, whose product is at most $d(d+1)m^{2H}B^{2L}$. If one is zero, use the zero implementation. Otherwise rescale the layers by positive factors with product one to make all energies equal to their geometric mean. This preserves both the function and its active support. The amplitude bound follows directly from the original coefficient-bounded teacher.
\end{proof}

\begin{lemma}[Selective perturbation of a sparse mean network]
\label{ag-lem:relu-profile-comparison}
Let $N/\omega\le1$ and let a deterministic ReLU network $w_*$
have, at layer $\ell$, at most $m_\ell\le N$ active coordinates.
Assume its first-layer mean is supported on the active rows, each
internal mean matrix is supported on active current rows and active
preceding columns, and its readout is supported on the final active
coordinates. Suppose every weighted layer energy is at most $e$.
For $0<r\le1$, retain all prior Gaussian variances except those
entering active hidden rows and all readout variances, whose standard
deviations are multiplied by $r$. Center the resulting product Gaussian
law $\nu_r$ at $w_*$. Then, with $L=H+1$,
\begin{align*}
 (2/\omega)\operatorname{KL}(\nu_r\Vert\mu)
 &\le Q_N(w_*)+2\left(d+2+\sum_{\ell=2}^Hm_\ell\right)\log(1/r),\\
 \sup_{x\in[-1,1]^d}
 \left(\mathbb E_{\nu_r}|z_w(x)-z_{w_*}(x)|^2\right)^{1/2}
 &\le\sqrt2 Lr(e+1)^{H/2}.
\end{align*}
\end{lemma}
\begin{proof}
There are
$k=m_1(d+1)+N\sum_{\ell=2}^Hm_\ell+N$ narrowed coordinates.
The exact Gaussian entropy formula and $N/\omega\le1$ give
\[
 (2/\omega)\operatorname{KL}(\nu_r\Vert\mu)
 =Q_N(w_*)+(1/\omega)k(r^2-1-2\log r)
 \le Q_N(w_*)+2\left(d+2+\sum_{\ell=2}^Hm_\ell\right)\log(1/r).
\]
Fix an input and set $S_\ell=\|h^\ell\|_2^2/N$. The normalized
first mean-layer squared norm is at most $2e$; its Gaussian noise
contributes at most two to its expected squared norm. Centering of
all Gaussian noise and $\operatorname{ReLU}(t)^2\le t^2$ imply
$\mathbb E S_1\le2(e+1)$. Each later mean matrix has operator norm
at most $\sqrt e$, and its conditional noise contributes at most
$(N/\omega)S_{\ell-1}\le S_{\ell-1}$. Hence
$\mathbb E S_\ell\le2(e+1)^\ell$.

Let $P_\ell$ project onto that layer's active coordinates, and put
$\Delta_\ell=(\mathbb E\|P_\ell h^\ell-h_*^\ell\|_2^2/N)^{1/2}$.
The deterministic hidden state is supported on the active coordinates.
Lipschitzness and conditional independence give
\[
 \Delta_1\le\sqrt2r,\qquad
 \Delta_\ell\le\sqrt e\,\Delta_{\ell-1}
       +r\sqrt{m_\ell/\omega}\sqrt{\mathbb E S_{\ell-1}}
 \le\sqrt e\,\Delta_{\ell-1}
       +\sqrt2r(e+1)^{(\ell-1)/2}.
\]
Thus $\Delta_\ell\le\sqrt2\ell r(e+1)^{(\ell-1)/2}$.
The supported mean readout has $\sqrt N\|v_*\|_2\le\sqrt e$,
while its independent perturbation contributes output variance
$(r^2/\omega)S_H$. Consequently the total root mean square output
error is at most
$\sqrt e\,\Delta_H+r\sqrt{(1/\omega)\mathbb E S_H}
\le\sqrt2 Lr(e+1)^{H/2}$.
\end{proof}

\begin{corollary}[Homogeneous resource dependence for every $\gamma\ge1$]
\label{ag-cor:relu-cost}
For ReLU, Theorem~\ref{ag-thm:main} holds with the following sharper comparison cost in place of $A_{m,b,n}$:
\[
 A^{\rm R}_{m,b,n}=1+E_{\rm R}
 +2[d+2+(H-1)m]\log[4K_e(U_{\rm R}+1)n],
 \qquad K_e=\sqrt2L(e+1)^{H/2},
\]
where $B=b$ in the teacher quantities. The same replacement is made in $\mathcal B_{m,b,n}$. Its remainder is
\[
 \widetilde O\left(
 m^{2H/L}b^2 R^{1-1/L}n^{-2/[L(d+2)]}
       +m^{2H}b^{2L}n^{-2/(d+2)}\right)
\]
at fixed dimension and depth. In particular the noiseless resource power is $m^{2H}b^{2L}$, and at positive fixed risk the growing-cap sufficient range is $H\alpha+L\beta<1/(d+2)$ for $m\le n^\alpha,b\le n^\beta$.
\end{corollary}
\begin{proof}
Apply Lemma~\ref{ag-lem:relu-profile-comparison} to the balanced teacher, with all active sizes equal to $m$, and choose $r=1/[4K_e(U_{\rm R}+1)n]$. Its pointwise root mean square output error is at most $1/[4(U_{\rm R}+1)n]$. Expanding square loss gives comparison expected empirical loss at most $\widehat R_S(g)+1/n$. The relative-entropy bound and \eqref{ag-eq:variational} therefore give $F_S(\lambda)\le\widehat R_S(g)+\lambda A^{\rm R}_{m,b,n}$ whenever $\lambda\ge1/n$. The selected root meets this condition. Every subsequent step of Theorem~\ref{ag-thm:main} uses only this comparison and $(U_{\rm R}+1)^2\le C E_{\rm R}^L$, so it applies unchanged with $r=L$. Finally $2H/L\ge1$, hence the logarithmic covariance cost is of order at most $\widetilde O(E_{\rm R})$, proving the displayed powers.
\end{proof}


\section{One posterior can use the best available certificate}
\label{sec:certificate-selection}
The two arguments apply to the same Gaussian prior and likelihood family. At a scale where both are valid, evaluate each calibration with failure probability $\delta/2$. The full-energy certificate is
\[
 C_E(S)=F_S(\lambda_E)+c_0\lambda_E+n^{-4}
\]
for $r\ge2$, or
\[
 C_E(S)=F_S(\lambda_E)(1+4a_n^2/\lambda_E)+c_0\lambda_E+n^{-4}
\]
for shallow tanh. The operator-norm certificate is
\[
 C_O(S)=F_S(\lambda_O)+2B_H\lambda_O+n^{-3},
\]
or its declared deterministic bound on the fallback branch. All quantities are functions of the specified learner and the observed sample.

Select whichever pair $(\lambda_E,C_E)$ or $(\lambda_O,C_O)$ has smaller certificate, with a fixed tie rule. On the intersection of the two theorem events, the posterior-average risk at the selected temperature is at most $\min\{C_E,C_O\}$. Its ordinary mean has no larger risk by Jensen. The comparison arguments bound the certificates themselves, so the selected learner inherits the smaller of the corresponding approximation-oracle bounds. The union of their failure probabilities is at most $\delta$.

If only the operator-norm argument applies, use that certificate alone. No teacher cap, risk, or noise parameter is supplied to this selection rule, and no temperature mixture is introduced. For example, a deep mean-field network uses the all-scale argument; a sufficiently strong tanh prior can additionally use the sharper full-energy certificate.

\section{Interpretation and scope}
The statistical gain comes from the actual posterior objective. When a small comparison network has small risk, its comparison measure gives small free energy. The weighted objective-tail bound then reduces both the empirical-loss range and the parameter range that need to be controlled. The calibration adjusts to that reduction using only the observed data.

For the full-energy argument an exactly fitting teacher gives a $1/n$ rate, while the all-scale argument gives $n^{-2/(d+2)}$. At positive fixed comparison risk, their slower depth-dependent exponents remain. These statements concern risk, so they also cover approximation error and bounded label noise without assuming that the Gaussian likelihood is a correct data-generating model.

The coefficient cap matters for unrestricted input laws. A fixed number of neurons with very large coefficients can vary on very small input scales. The present theorems do not identify network width alone with an adequate uniform complexity measure. Separate bounded-output results under regular input laws can change that conclusion; those are different observation conventions and assumptions.

The prior is fixed at each learner architecture. The likelihood variance depends on sample size and is selected from free energy for prediction. It is not claimed to estimate the physical conditional noise variance. Appendix~\ref{sec:normalization} gives simple obstructions to several fixed-variance and fixed-hyperprior alternatives.

\paragraph{Verification status.}
These are conventional mathematical proofs, not Lean formalizations. The two main statistical routes and the activation extension have undergone separate mathematical audits. The full-energy route invokes the stated smooth-loss Rademacher inequality; the operator-norm route constructs its statistical event directly from finite covers and Bernstein's inequality. All events are fixed before the sample-dependent likelihood calibration.

\appendix
\section{A sharp path-variation specialization for one hidden ReLU layer}
\label{sec:shallow-ReLU-path}
For one hidden layer, the original independent Gaussian prior already gives a quadratic path-variation oracle without a soft gain penalty. Inputs lie in the Euclidean unit ball of $\R^d$, labels are in $[-1,1]$, and
\[
 z_w(x)=\sum_{j=1}^N v_j\operatorname{ReLU}(w_j\cdot x),\qquad
 w_{ji}\sim\mathcal N(0,\tau^2/d),\quad
 v_j\sim\mathcal N(0,\tau^2/N^2),\quad0<\tau\le1.
\]
All coordinates are independent. A constant input coordinate can encode a first-layer bias after normalizing the augmented input ball. For a comparison implementation $g$, write
\[
 P=\sum_{j=1}^C|v_j|\|w_j\|_2.
\]

\begin{corollary}[Quadratic path variation with additive input dimension]
There is a free-energy calibration, independent of the comparison width, coefficients, and risk, such that the ordinary raw posterior mean has
\[
 R(\widehat f)\le R(g)+\widetilde O\left(
 (P+\sqrt d+1)\sqrt{R(g)/n}
 +(P+\sqrt d+1)^2/n\right)
\]
with probability at least $1-\delta$, simultaneously for every $N\ge C$ and every $0<\tau\le1$. Thus a sufficient sample count for excess $\varepsilon$ is
\[
 \widetilde O\left((P^2+d+1)
             \{R(g)/\varepsilon^2+1/\varepsilon\}\right).
\]
Only logarithms in $n,d,P,1/\varepsilon,1/\delta$ are suppressed. In particular, bounded first-row norms and readout coefficients give quadratic dependence on comparison width. The Gaussian prior remains unchanged.
\end{corollary}
\begin{proof}
Use the energy and precision normalization
\[
 u=\frac dN\|W\|_F^2,\qquad q=N\|v\|^2,
 \qquad Q=\frac{u+q}{\sqrt d},\qquad
 \omega=\frac{N\sqrt d}{\tau^2}.
\]
Then $\omega Q$ has a chi-square law with $N(d+1)$ degrees of freedom. Its normalized degrees of freedom are at most
\[
 D_0=(d+1)/\sqrt d.
\]
Moreover $P(w)\le\|W\|_F\|v\|=\sqrt{uq/d}\le Q/2$, so $|z_w(x)|\le Q/2$. The one-layer path class has worst-case Rademacher complexity at most $K_1Q/\sqrt n$ for an absolute $K_1$: ReLU contraction and absolute convexification reduce to Euclidean-unit linear features. The centered independent readout also gives prior population and empirical risk at most two.

Positive rescaling of each teacher neuron makes its two weighted energy contributions equal, without changing the function. At every $N\ge C$, zero-pad the remaining neurons. The resulting first-layer and readout means satisfy
\[
 u=q=\sqrt dP,\qquad Q=2P.
\]
Use these means and multiply every native coordinate variance by $\eta^2$, for $0<\eta<1$. Its normalized Gaussian relative entropy is bounded by
\[
 \frac2\omega\operatorname{KL}(\rho\|\mu)
 =2P+\frac{N(d+1)}\omega(\eta^2-1-\log\eta^2)
 \le2P+2D_0\log(1/\eta).
\]
For the uniform perturbation error, first-layer noise has expected squared Frobenius norm $\eta^2\tau^2N$, and readout noise has expected squared norm $\eta^2\tau^2/N$. Independence, ReLU Lipschitzness, and the displayed mean energies give
\[
 \left(\E_\rho\|z_w-g\|_{\infty}^2\right)^{1/2}
 \le\eta\tau d^{1/4}\sqrt P
       +\eta\tau d^{-1/4}\sqrt P+\eta^2\tau^2
 \le2\eta d^{1/4}\sqrt P+\eta^2.
\]
Take $\eta=[64n(d+1)(P+1)^2]^{-1}$. Since $|g|\le P$, the empirical squared-risk increase is at most $1/(2n)$. Therefore, with Gaussian observation variance $n\lambda/\omega$,
\[
 F_S(\lambda)\le\widehat R_S(g)+1/(2n)+\lambda T_g,
 \qquad T_g=6+2P+2D_0\log(1/\eta).
\]

Apply the chi-square-energy adaptive certificate with capacity radius linear in $Q$ (the $r=2$ case). Its root is
\[
 \widehat\lambda^2
 =256a_n^2[2F_S(\widehat\lambda)+\widehat\lambda(K_n+1)],
\]
where $a_n=\widetilde O(n^{-1/2})$, $K_n=O(D_0+\log(n(D_0+1)))$, and the certificate's additive coefficient is $c_0=1/8+D_0\log2/4$. Its empirical comparison budget is
\[
 \mathcal B_g=2T_g+K_n+1+c_0
             =\widetilde O(P+\sqrt d+1).
\]
The weighted-objective tail proof uses only $\omega Q\sim\chi^2$ and $N(d+1)/\omega\le D_0$; it does not require $D_0\le1$. The smooth-loss event and fixed-comparator Bernstein step therefore give the stated population bound. These events are common to all widths and scales because the path-capacity class already contains all widths, and the displayed bounds use only $\tau\le1$. Jensen transfers posterior-average risk to the ordinary mean.
\end{proof}

\section{Sub-Gaussian labels without changing the parameter prior}
\label{sec:subgaussian-adaptive-transfer}
The raw-output adaptive Gaussian oracle extends to unbounded labels with a known marginal sub-Gaussian tail. Its powers of sample size and teacher complexity are unchanged; only logarithmic factors and an explicit negligible label-tail term are added. The learner still trains on the original observations and predicts the ordinary raw posterior mean. No Gaussian weight or prior is rescaled.

We state the transfer for the full-energy and input-covering theorems proved here, and for finite-program variants with the same statistical structure. The common-Lipschitz and bounded-smooth-activation finite-program theorems of~\cite{common-program,smooth-program} also satisfy its hypotheses. Suppose the bounded-label theorem has powers $v,\rho>0$, $r=v/\rho\ge2$, deterministic statistical factor $\theta_n$, and a deterministic comparison resource $\bar A$ such that its posterior-average oracle is
\begin{equation}
 G(S):=\int R(z_w)d\Pi_S
 \le R(g)+C\left[
 \bar A R(g)^{1-1/r}\theta_n^{1/v}
          +\bar A^r\theta_n^{1/\rho}\right]+n^{-3}.
 \label{eq:transfer-unit-oracle}
\end{equation}
The theorem's confidence parameter, including all scalar-comparator events, will be $\delta/2$. Its fixed prior family, admissible learner widths/scales, and teacher class are unchanged. The needed closure property is that its proof remains valid for the observation map $w\mapsto z_w/T$ and teacher $g/T$, using the same original parameter vectors and comparison resource, whenever $T\ge1$. We verify that property below for the Gaussian dictionary and full-energy arguments.

Assume that for a known $s_Y>0$,
\begin{equation}
 \Pr(|Y|>t)\le2e^{-t^2/(2s_Y^2)}\quad(t\ge0).
 \label{eq:transfer-label-tail}
\end{equation}
Let $n\ge2$ and set
\begin{equation}
 T=\max\{1,s_Y\sqrt{8\log(8n/\delta)}\},
 \qquad u=2s_Y\left(\frac\delta{8n}\right)^2.
 \label{eq:transfer-T}
\end{equation}
For a physical likelihood parameter $\lambda$, write $F_S(\lambda)$ for its raw free energy. Run the bounded-label calibration on the normalized observable and labels, namely
\[
 z'_w=z_w/T,\qquad Y'_i=Y_i/T,
 \qquad F'_S(\lambda')=T^{-2}F_S(T^2\lambda').
\]
If it selects $\lambda'$, the actual learner uses $\lambda=T^2\lambda'$. This is an observable scalar likelihood calibration. In particular, for the dictionary theorem the root is applied to $F'_S$, not to an unmodified free energy with silently changed units. The physical Gaussian observation variance is $nsT^2\lambda'$ under prior scale $s$.

For $x\ge0$, define
\begin{equation}
 e_T(x)=C\left[
 \bar A T^{2/r}x^{1-1/r}\theta_n^{1/v}
       +\bar A^rT^2\theta_n^{1/\rho}\right]+T^2n^{-3}.
 \label{eq:transfer-remainder}
\end{equation}
\begin{theorem}[Raw adaptive oracle with sub-Gaussian labels]
\label{thm:subgaussian-adaptive-transfer}
Under the stated bounded-label theorem and its verified observation-map closure, the posterior trained on the original unbounded labels satisfies, with probability at least $1-\delta$,
\begin{equation}
 \int R(z_w)d\Pi_S
 \le R(g)+e+4u\sqrt{R(g)+e}+4u^2,
 \qquad e=e_T\big((\sqrt{R(g)}+u)^2\big).
 \label{eq:transfer-oracle}
\end{equation}
The ordinary raw posterior mean has the same risk bound. If the bounded-label event is common to all admissible widths and scales, so is this event. At fixed label-tail parameter $s_Y$, depth and confidence, the powers $1/v$ and $1/\rho$, and every polynomial teacher-resource exponent, are unchanged up to logarithms.
\end{theorem}
\begin{proof}
Let $Y_T=\operatorname{clip}_{[-T,T]}Y$ and let $R_T$ be its population squared risk. Tail integration gives
\[
 \|Y-Y_T\|_2^2
 =\int_0^\infty2a\Pr(|Y|>T+a)\,da
 \le4s_Y^2e^{-T^2/(2s_Y^2)}\le u^2.
\]
Also a union bound gives
\[
 \Pr(\exists i:Y_i\ne(Y_i)_T)
 \le2ne^{-T^2/(2s_Y^2)}
 \le2n(\delta/(8n))^4<\delta/2.
\]
On the complementary event the original and clipped training samples coincide exactly, so every prior, likelihood, free energy, selected temperature and posterior coincides as well. Truncation of labels is used only in the proof.

Apply the bounded-label theorem to $(X,Y_T/T)$ and observable $z_w/T$. Its posterior likelihood is exactly the actual likelihood at $\lambda=T^2\lambda'$, since
\[
 \frac{\widehat R_S(z_w)}{2\lambda s}
 =\frac{\widehat R_{S/T}(z_w/T)}{2\lambda' s}.
\]
Multiplying \eqref{eq:transfer-unit-oracle} by $T^2$ yields
\[
 \int R_T(z_w)d\Pi_S\le R_T(g)+e_T(R_T(g)).
\]
Minkowski gives $\sqrt{R_T(g)}\le\sqrt{R(g)}+u$.
The function $e_T$ is nondecreasing, so the right side is at most
$R_T(g)+e$, with $e$ as in the theorem.

For each observed sample, apply Minkowski on the product of the posterior law and an independent test observation:
\[
 \sqrt{\int R(z_w)d\Pi_S}
 \le\sqrt{\int R_T(z_w)d\Pi_S}+u
 \le\sqrt{R(g)+e}+2u.
\]
The last step uses $\sqrt{R_T(g)+e}\le\sqrt{R(g)+e}+u$.
Squaring proves \eqref{eq:transfer-oracle}. The two exceptional events total at most $\delta$. Jensen proves the ordinary-mean statement. All posterior moments invoked in the raw Gaussian theorem remain finite; the comparison of the original and clipped posteriors is exact on the stated event.
\end{proof}

\subsection{Why the new Gaussian proofs have the required closure}
The normalization acts only on the observation map, not on parameters. Each original network and teacher is retained, and its output is divided by $T$ in the analysis.

For the finite-program proof, divide every dictionary center by $T$. Its cardinality does not change, its amplitude decreases, and an original $L^2$ approximation error at most $a$ becomes at most $a/T\le a$. The prior-good-set probabilities, normalized parameter energies and prior-tail bounds are unchanged. The normalized-observable prior risk under labels bounded by one is no greater than the same constant used in the bounded-label proof, after retaining the constant label contribution. Consequently the fixed-dictionary statistical event, sample-dependent complexity, weighted objective tail, and scalar free-energy root apply with the same $p,q,v,\rho$ and the same upper statistical constants. This includes the improved bounded-$C^2$ dictionary: its smaller entropy exponents are unaffected by dividing outputs by $T$.

The Gaussian comparison also keeps the original teacher parameters. Its KL cost is unchanged. More explicitly, suppose the native comparison has output RMS error at most $e_0$, with $|g|\le U$, and the bounded-label proof chose it so that
$2(U+1)e_0+e_0^2\le1/n$. The normalized squared-loss increment is at most
\[
 2(U/T+1)e_0/T+e_0^2/T^2
 \le2(U+1)e_0+e_0^2\le1/n.
\]
Thus the same comparison resource $\bar A$ is sufficient for the normalized theorem; no new teacher cap or width factor is introduced. The teacher-amplitude condition also improves under division by $T$.

For the full-energy proof, empirical Rademacher complexity, amplitude and prior raw moments decrease under output division by $T$. Its Gaussian energy identity and entropy cost are identical. The same argument therefore applies. The sharper shallow-tanh oracle transforms to
\[
 R_T(g)+\widetilde O\left(
 T\sqrt{\bar A R_T(g)/n}+T^2\bar A/n\right),
\]
followed by the same $L^2$ label-tail transfer. Its sample powers also remain unchanged.

The tail assumption includes Gaussian noise around a uniformly bounded conditional mean, with a sufficiently large known $s_Y$. No correctly specified observation-noise variance is required. The actual prior and Gaussian likelihood form are unchanged, but their selected likelihood variance acquires the explicit factor $T^2$ and the corresponding normalized free-energy calibration.

\section{Why the likelihood scale must depend on the prior}
\label{sec:normalization}
This appendix concerns bounded hidden activations, including tanh. It explains why the positive theorem uses a deliberately calibrated observation variance. The statements require no assumptions on the hidden-weight prior beyond independence from the centered Gaussian readout.

Suppose the last hidden vector satisfies $|h_j(x)|\le1$ and the readout prior is
$v\sim\mathcal N(0,(s_N/N)I_N)$, with $s_N\to0$. Conditional on hidden weights, the output at one input is a centered Gaussian of variance at most $s_N$.

\begin{proposition}[Fixed observation variance cannot give an all-width oracle]
\label{prop:fixed-variance-failure}
Train on $n$ copies of an input $x_0$ with label one, using a Gaussian observation variance $V>0$. At every finite width, the ordinary posterior mean at $x_0$ lies in
\[
 \left[0,\frac{ns_N}{V+ns_N}\right].
\]
Its test loss on the same noiseless law is therefore at least
\begin{equation}
 \left(\frac{V}{V+ns_N}\right)^2.
 \label{eq:variance-necessary}
\end{equation}
In particular, vanishing loss along such a sequence requires $V/(ns_N)\to0$. For the positive theorem's variance $V=n\lambda s_N$, this lower bound is $(\lambda/(1+\lambda))^2$.
\end{proposition}
\begin{proof}
Conditional on hidden weights, put $q=(s_N/N)\|h(x_0)\|^2\le s_N$. Gaussian conjugacy gives conditional posterior output mean $nq/(V+nq)$. Integrating the posterior law of hidden weights preserves the stated interval.
\end{proof}

A width-one tanh teacher with coefficients at most two realizes the constant-one target at every depth: keep each hidden value at $1/2$ using first bias $\operatorname{arctanh}(1/2)$ and subsequent self-edge $\log3$, and use readout weight two. Thus \eqref{eq:variance-necessary} already concerns a uniformly simple teacher.

The repeated-input example gives an exact bound, but the basic failure does not depend on a degenerate input law. On an arbitrary sample with $|Y_i|\le1$, conditional Gaussian readout regression has mean
\[
 k_x^{\mathsf T}(K+VI_n)^{-1}y,
\]
where $K$ is the prior output covariance matrix on training inputs and every component of $k_x$ has absolute value at most $s_N$. Its absolute value is at most $ns_N/V$, uniformly in $x$ and hidden weights. Therefore the raw posterior mean tends uniformly to zero whenever $V/(ns_N)\to\infty$. A constant-one target can consequently be used with any input distribution, including a nondegenerate continuous one. Taking a fractional power $0<\beta\le1$ of a fixed Gaussian likelihood replaces $V$ by $V/\beta$ and cannot remove this obstruction.

\begin{proposition}[A fixed physical variance hyperprior]
\label{prop:variance-hyperprior}
Put any proper probability hyperprior $\pi(dV)$ on $(0,\infty)$, independent of learner width and weights, and use the fully normalized Gaussian likelihood. Fix $n\ge2$ repeated-input observations at $x_0$, with label mean $\bar y$ and positive residual sum of squares
\[
 A=\sum_i(y_i-\bar y)^2>0.
\]
Then the ordinary posterior output mean at $x_0$ tends to zero as $s_N\to0$, regardless of the hidden priors.
\end{proposition}
\begin{proof}
After integrating the Gaussian readout conditional on hidden weights, the marginal likelihood is
\[
 L(q,V)=(2\pi)^{-n/2}V^{-(n-1)/2}(V+nq)^{-1/2}
 \exp\left[-\frac A{2V}-\frac{n\bar y^2}{2(V+nq)}\right],
 \qquad0\le q\le s_N.
\]
The corresponding conditional output mean is $m(q,V)=nq\bar y/(V+nq)$. For each fixed $V>0$, uniformly over $0\le q\le s_N$, one has $L(q,V)\to L(0,V)$ and $m(q,V)\to0$. Furthermore,
\[
 L(q,V)\le(2\pi)^{-n/2}V^{-n/2}e^{-A/(2V)},
 \qquad |m(q,V)|\le|\bar y|.
\]
The likelihood upper bound is bounded on $(0,\infty)$ because $A>0$, hence integrable against every proper $\pi$. Dominated convergence, uniformly over the hidden-prior mixtures, gives a positive limiting denominator $\int L(0,V)\pi(dV)$ and a numerator tending to zero. Their ratio is the posterior mean.
\end{proof}

For the correctly specified regression law $X=x_0$, $Y=1+\varepsilon$, with $\varepsilon\sim\mathcal N(0,a^2)$ and $a>0$, the residual $A$ is positive almost surely for every fixed $n\ge2$. Proposition~\ref{prop:variance-hyperprior} gives test excess tending to one over the constant-one teacher, almost surely over that sample. This example allows unbounded Gaussian labels only in the present appendix. Bounded noise $\varepsilon=\pm a$ gives the same conclusion on the event that both signs occur, of probability $1-2^{1-n}$.

\begin{proposition}[An unbounded hyperprior on normalized temperature]
\label{prop:temperature-hyperprior}
Instead give $\lambda$ any proper hyperprior $\pi_n$ on $(0,\infty)$ with unbounded support, independent of learner width, and set the physical observation variance to $n\lambda s_N$. The hyperprior may depend on sample count. Use the fully normalized Gaussian likelihood. For every fixed repeated-input sample as in Proposition~\ref{prop:variance-hyperprior}, the posterior temperature tends to infinity in probability, and the ordinary output mean tends to zero, as $s_N\to0$.
\end{proposition}
\begin{proof}
Write $t=q/s_N\in[0,1]$ and $M_2=n^{-1}\sum_i y_i^2=A/n+\bar y^2$. Apart from the factor $(2\pi ns_N)^{-n/2}$, the integrated likelihood is
\[
 \mathcal L_{s_N}(t,\lambda)=\lambda^{-(n-1)/2}(\lambda+t)^{-1/2}
 \exp\left[-\frac1{2s_N}\left(\frac A{n\lambda}+
                                      \frac{\bar y^2}{\lambda+t}\right)\right].
\]
Fix $K>0$. When $s_N\le A/(n^2K)$, differentiation shows that
$\lambda^{-n/2}e^{-A/(2ns_N\lambda)}$ is increasing on $(0,K]$. Hence, uniformly in $t$,
\[
 \sup_{\lambda\le K}\mathcal L_{s_N}(t,\lambda)
 \le K^{-n/2}e^{-A/(2nKs_N)}.
\]
Choose $L>2nKM_2/A$. Unbounded support supplies a finite $U>L$ with $\pi_n([L,U])>0$. Throughout this interval,
\[
 \mathcal L_{s_N}(t,\lambda)
 \ge U^{-(n-1)/2}(U+1)^{-1/2}e^{-M_2/(2Ls_N)}.
\]
The ratio bounds posterior probability $\Prob\{\lambda\le K\mid S\}$ by a finite constant times $e^{-A/(4nKs_N)}$. This tends to zero. Conditional readout mean has magnitude
\[
 |\bar y|\frac{t}{\lambda+t}\le\frac{|\bar y|}{1+\lambda}.
\]
First split its posterior expectation at $K$, then let $s_N\to0$ and $K\to\infty$. The mean converges to zero.
\end{proof}

These propositions concern ordinary variance or temperature hyperpriors with the stated width independence. They do not cover a width-dependent hyperprior, a temperature prior with a shrinking upper endpoint, or a prediction algorithm that weights candidate temperatures using a separate validation criterion. The free-energy calibration in the main theorem is different: it explicitly balances a statistical error certificate and selects one posterior. The distinction is necessary even for constant regression.

\begin{thebibliography}{99}
\bibitem{sst} Nathan Srebro, Karthik Sridharan and Ambuj Tewari.
\emph{Smoothness, Low Noise and Fast Rates}. Advances in Neural Information Processing Systems 23, 2010.
\url{https://www.cs.cornell.edu/~sridharan/smoothness.pdf}.
\bibitem{talagrand} Michel Talagrand. \emph{Upper and Lower Bounds for Stochastic Processes}. Springer, 2014. Gaussian Sudakov minoration, entropy bounds, and the two-sided majorizing-measure theorem.
\bibitem{finite} ChatGPT-6 Astra, Claude Opus 5.5, Kaarel H\"anni and Dmitry Vaintrob.
\emph{Finite Gaussian Bayesian networks can compete with small approximating networks}. Research manuscript, 2026.
\url{https://github.com/kaarelh/AI-works/tree/main/neural-network-generalization/works-in-progress/agnostic-finite-posteriors}.
\bibitem{relu} ChatGPT-6 Astra, Claude Opus 5.5, Kaarel H\"anni and Dmitry Vaintrob.
\emph{A sharp comparator-growth threshold for Gaussian Bayesian ReLU networks}. Research manuscript, 2026.
\url{https://github.com/kaarelh/AI-works/tree/main/neural-network-generalization/works-in-progress/agnostic-relu-sharp}.
\bibitem{common-program} ChatGPT-6 Astra, Claude Opus 5.5, Kaarel H\"anni and Dmitry Vaintrob.
\emph{Agnostic oracle bounds for Gaussian neural networks with Lipschitz activations}. Research manuscript, 2026.
\bibitem{smooth-program} ChatGPT-6 Astra, Claude Opus 5.5, Kaarel H\"anni and Dmitry Vaintrob.
\emph{Finite Bayesian neural networks as agnostic approximation learners}. Research manuscript, 2026.
\end{thebibliography}
\end{document}
